\documentclass[11pt,twocolumn,a4paper]{article}

%% ── Packages ──────────────────────────────────────────────────────────────────
\usepackage[utf8]{inputenc}
\usepackage[T1]{fontenc}
\usepackage{lmodern}
\usepackage{microtype}
\usepackage{amsmath,amssymb,amsthm,mathtools}
\usepackage{physics}
\usepackage{bm}
\usepackage{graphicx}
\usepackage{xcolor}
\usepackage{booktabs}
\usepackage{multirow}
\usepackage{enumitem}
\usepackage{float}
\usepackage{algorithm}
\usepackage{algorithmic}
\usepackage[margin=1.8cm,columnsep=0.6cm]{geometry}
\usepackage[colorlinks=true,linkcolor=blue!60!black,citecolor=green!50!black,urlcolor=blue!70!black]{hyperref}
\usepackage{natbib}
\usepackage{authblk}
\usepackage{fancyhdr}
\usepackage{titlesec}
\usepackage{caption}
\usepackage{subcaption}

%% ── Theorem environments ──────────────────────────────────────────────────────
\newtheorem{axiom}{Axiom}
\newtheorem{theorem}{Theorem}
\newtheorem{lemma}[theorem]{Lemma}
\newtheorem{corollary}[theorem]{Corollary}
\newtheorem{proposition}[theorem]{Proposition}
\newtheorem{definition}{Definition}
\newtheorem{remark}{Remark}

%% ── Custom commands ───────────────────────────────────────────────────────────
\newcommand{\Ocat}{\hat{O}_{\mathrm{cat}}}
\newcommand{\Ophys}{\hat{O}_{\mathrm{phys}}}
\newcommand{\Sigvis}{\Sigma_{\mathrm{vis}}}
\newcommand{\Siginv}{\Sigma_{\mathrm{inv}}}
\newcommand{\dcat}{d_{\mathrm{cat}}}
\newcommand{\Sk}{S_k}
\newcommand{\St}{S_t}
\newcommand{\Se}{S_e}
\newcommand{\Stot}{S_{\mathrm{total}}}
\newcommand{\R}{\mathbb{R}}
\newcommand{\N}{\mathbb{N}}
\newcommand{\Z}{\mathbb{Z}}
\newcommand{\hilbert}{\mathcal{H}}
\newcommand{\partspace}{\mathcal{P}}
\newcommand{\catspace}{\mathcal{C}}
\newcommand{\phasespace}{\Gamma}

%% ── Title information ─────────────────────────────────────────────────────────
\title{\textbf{Measurement-Modality Stereograms: Dual-Path Pixel Validation Through Optical and Oxygen-Mediated Categorical Observation}}

\author{
Kundai Farai Sachikonye\\
AIMe Registry for Artificial Intelligence\\
\texttt{kundai.sachikonye@bitspark.com}
}

\date{\today}

%% ── Header/footer ─────────────────────────────────────────────────────────────
\pagestyle{fancy}
\fancyhf{}
\fancyhead[L]{\small\textit{Measurement-Modality Stereograms}}
\fancyhead[R]{\small\thepage}
\renewcommand{\headrulewidth}{0.4pt}

%% ══════════════════════════════════════════════════════════════════════════════
\begin{document}
\maketitle
\thispagestyle{fancy}

%% ── Abstract ──────────────────────────────────────────────────────────────────
\begin{abstract}
Every pixel in a biological microscopy image carries an implicit duality: the value recorded by the camera sensor through external photon collection (the \emph{visible pixel}) and the value that the intracellular oxygen distribution at that same spatial location would encode through molecular state transitions (the \emph{invisible pixel}).
We develop a rigorous mathematical framework---\emph{measurement-modality stereograms}---that formalises this duality from two foundational axioms: bounded phase space and categorical observation.
From these axioms we derive hierarchical partition coordinates $(n, \ell, m, s)$ for spherical phase space with cumulative capacity $C(n)=2n^2$, a conserved entropy coordinate system $(\Sk, \St, \Se) \in [0,1]^3$, and a commutation theorem establishing that categorical observables commute with physical observables: $[\Ocat, \Ophys] = 0$.
We prove that molecular oxygen, possessing two relevant electronic levels, admits exactly three categorical states---absorption (detector), ground (phase reference), and emission (virtual source)---forming a complete ternary basis with information content $I = N\log_2 3$ per $N$ molecules.
The visible pixel is defined by external optical detection, and the invisible pixel by oxygen-mediated mid-infrared emission at $\lambda \approx 3.0\;\mu$m.
Both independently encode the same partition signature, enabling dual-pixel cross-validation.
We prove a consistency theorem requiring agreement between the two paths for all reliable measurements, an information gain theorem quantifying the strictly greater information content of the fused dual pixel, and a resolution enhancement theorem exploiting the sub-diffraction scale of molecular oxygen detectors (${\sim}0.1$~nm) relative to the optical diffraction limit (${\sim}200$~nm).
Finally, we develop the stereogram fusion framework, demonstrate that disagreement patterns between visible and invisible pixels encode three-dimensional structural information analogous to binocular disparity, and propose a validation protocol using the BBBC039 nuclei segmentation dataset.

\medskip
\noindent\textbf{Keywords:} dual-pixel validation, categorical observation, oxygen sensing, partition coordinates, sub-diffraction imaging, measurement stereogram, information gain, microscopy, cross-modal fusion
\end{abstract}

%% ══════════════════════════════════════════════════════════════════════════════
\section{Introduction}\label{sec:introduction}
%% ══════════════════════════════════════════════════════════════════════════════

Modern biological microscopy produces images whose individual pixel values are accepted almost axiomatically as faithful reports of sample structure.
A brightfield pixel records the attenuation of transmitted light; a fluorescence pixel records emitted photon counts after excitation; a phase-contrast pixel records phase shifts induced by refractive index variations.
In every case, the measurement pathway is \emph{optical and external}: photons originate from a source, interact with the specimen, and are collected by a detector placed outside the sample \citep{Murphy2012,Mertz2019}.

Yet each spatial location within a biological specimen is simultaneously the site of continuous molecular activity.
Dissolved oxygen molecules undergo electronic state transitions, absorb and emit photons in the mid-infrared, and redistribute themselves according to local metabolic demand \citep{Wilson2008,Carreau2011}.
This molecular activity constitutes an independent \emph{internal} measurement channel---one that encodes spatial structure through oxygen concentration gradients, molecular state populations, and intracellular radiative transfer at wavelengths far from the visible band.

The central insight of this paper is that these two measurement channels---the external optical channel and the internal oxygen-mediated channel---are not merely complementary data sources to be combined heuristically.
Rather, they are \emph{independent observers of the same underlying categorical partition}, and their agreement or disagreement carries rigorous information-theoretic content.
We term the pair a \emph{measurement-modality stereogram}: not a stereogram in the spatial-parallax sense of binocular vision, but one in a deeper measurement-theoretic sense, where ``depth'' emerges from the parallax between an outside-in and an inside-out observation of the same physical reality.

% ── FIGURE PROPOSAL 1 ──
% Figure 1: Conceptual overview. Left panel: standard microscopy (external photon path
% from illumination source -> specimen -> objective -> camera sensor, producing the
% "visible pixel"). Right panel: oxygen-mediated internal path (O2 molecules at the same
% spatial location undergoing ternary state transitions, encoding structure via mid-IR
% emission, producing the "invisible pixel"). Centre panel: fusion into measurement-
% modality stereogram with consistency check symbol. Arrows should emphasise the
% independence of the two physical pathways converging on the same partition coordinates.

To place this framework on rigorous footing, we build the entire theory from two axioms (Section~\ref{sec:axioms}): \emph{bounded phase space} and \emph{categorical observation}.
From these, we derive hierarchical partition coordinates (Section~\ref{sec:partition}), a conserved entropy coordinate system (Section~\ref{sec:entropy}), and the commutation relation between categorical and physical observables (Section~\ref{sec:commutation}).
We then develop the ternary molecular state basis for oxygen (Section~\ref{sec:ternary}), formally define the visible pixel (Section~\ref{sec:visible}) and the invisible pixel (Section~\ref{sec:invisible}), prove dual-pixel consistency (Section~\ref{sec:consistency}), information gain (Section~\ref{sec:information}), and resolution enhancement (Section~\ref{sec:resolution}) theorems.
We construct the stereogram framework (Section~\ref{sec:stereogram}), define the cross-modal interaction potential (Section~\ref{sec:interaction}), extract depth from dual-pixel disagreement (Section~\ref{sec:depth}), and propose a concrete validation protocol (Section~\ref{sec:validation}).

No external framework or prior publication by the author is assumed.
Every definition, theorem, and derivation proceeds from first principles within this paper.


%% ══════════════════════════════════════════════════════════════════════════════
\section{Foundational Axioms}\label{sec:axioms}
%% ══════════════════════════════════════════════════════════════════════════════

We begin with two axioms that are physically motivated and mathematically minimal.

\begin{axiom}[Bounded Phase Space]\label{ax:bounded}
Every physical system with finite total energy $E < \infty$ occupies a region $\phasespace_E \subset \phasespace$ of the full phase space $\phasespace$ such that the Liouville measure satisfies
\begin{equation}\label{eq:bounded}
\mu(\phasespace_E) = \int_{\phasespace_E} \prod_{i=1}^{f} dq_i\, dp_i < \infty,
\end{equation}
where $f$ is the number of degrees of freedom and $(q_i, p_i)$ are canonical coordinates.
\end{axiom}

This axiom is a direct consequence of the virial theorem for bound systems \citep{Goldstein2002} and of the Heisenberg uncertainty principle for quantum systems \citep{Heisenberg1927}.
A system with energy $E$ cannot explore arbitrarily large phase-space volumes because both position and momentum are constrained: positions are confined by the potential well, and momenta are bounded by $|p| \leq \sqrt{2mE}$.

\begin{axiom}[Categorical Observation]\label{ax:categorical}
An observer with finite spatial, temporal, and energetic resolution partitions the accessible phase space $\phasespace_E$ into a finite collection of distinguishable categories
\begin{equation}\label{eq:partition}
\mathcal{P} = \{C_1, C_2, \ldots, C_K\}, \quad \bigcup_{k=1}^{K} C_k = \phasespace_E, \quad C_i \cap C_j = \emptyset \;\; (i \neq j),
\end{equation}
and every measurement outcome is a category label $k \in \{1, 2, \ldots, K\}$.
\end{axiom}

Axiom~\ref{ax:categorical} formalises the fact that no physical detector has infinite precision.
Every measurement apparatus---from a CCD pixel to a molecular receptor---responds within finite dynamic range and with finite discrimination, so that distinct physical states mapping to the same detector response are observationally equivalent.
This echoes the coarse-graining foundational to statistical mechanics \citep{Boltzmann1877,Jaynes1957} and the finite-alphabet constraint in information theory \citep{Shannon1948}.

\begin{remark}
Axioms~\ref{ax:bounded} and~\ref{ax:categorical} together guarantee that $K < \infty$: a finite-measure phase space partitioned by finite-resolution bins yields finitely many categories.
This finiteness is essential for all information-theoretic results that follow.
\end{remark}

\subsection{Consequences of the axioms}

The two axioms have immediate consequences that structure the remainder of the paper.

\begin{lemma}[Finite information content]\label{lem:finite_info}
Any single observation of a system satisfying Axioms~\ref{ax:bounded} and~\ref{ax:categorical} carries at most
\begin{equation}\label{eq:maxinfo}
I_{\max} = \log_2 K \quad \text{bits}
\end{equation}
of information, where $K$ is the number of categories in the partition.
\end{lemma}
\begin{proof}
By Axiom~\ref{ax:categorical}, the observation outcome is an element of a finite set of cardinality $K$.
The maximum entropy of a discrete random variable on $K$ outcomes is $\log_2 K$ bits, achieved by the uniform distribution \citep{Shannon1948,CoverThomas2006}.
Since mutual information between the observation and the true state cannot exceed the entropy of the observation, $I \leq I_{\max} = \log_2 K$.
\end{proof}

\begin{lemma}[Observer-independence of category count]\label{lem:observer_indep}
The number of categories $K$ depends on (i) the total phase-space volume $\mu(\phasespace_E)$ and (ii) the observer's resolution volume $\delta\mu$, but not on the internal dynamics of the system:
\begin{equation}\label{eq:Kcount}
K = \left\lfloor \frac{\mu(\phasespace_E)}{\delta\mu} \right\rfloor.
\end{equation}
\end{lemma}
\begin{proof}
The partition in Axiom~\ref{ax:categorical} tiles the bounded phase space with non-overlapping cells of volume at least $\delta\mu$ (the observer's resolution limit).
The maximum number of such cells in a region of total volume $\mu(\phasespace_E)$ is $\lfloor \mu(\phasespace_E)/\delta\mu \rfloor$.
This count depends only on the ratio of total volume to resolution volume, not on the Hamiltonian governing the system's time evolution.
\end{proof}


%% ══════════════════════════════════════════════════════════════════════════════
\section{Partition Coordinates from Bounded Spherical Phase Space}\label{sec:partition}
%% ══════════════════════════════════════════════════════════════════════════════

We now derive a concrete coordinate system for the categorical partition when the bounded phase-space region has spherical symmetry.
This is the natural geometry for atomic, molecular, and cellular systems where a central potential or a cell boundary imposes approximate radial confinement \citep{Griffiths2018,Sakurai2020}.

\subsection{Hierarchical subdivision}

Consider a three-dimensional system with bounded, spherically symmetric phase space.
The natural coordinates are $(r, \theta, \varphi)$ in position space and $(p_r, p_\theta, p_\varphi)$ in momentum space.
By Axiom~\ref{ax:bounded}, both position and momentum have finite range.
We introduce a hierarchical subdivision defined by a principal level number $n \in \{1, 2, 3, \ldots\}$.

\begin{definition}[Principal level]\label{def:principal}
The principal level $n$ labels a radial shell in phase space containing all categories whose average phase-space radius falls within the $n$-th annular region.
The total phase-space volume up to level $n$ scales as $n^3$ for the position sector and is partitioned by angular structure in both sectors.
\end{definition}

At each principal level $n$, the angular structure in the position sector introduces an angular momentum quantum number:

\begin{definition}[Angular level]\label{def:angular}
For principal level $n$, the angular level $\ell$ takes values $\ell \in \{0, 1, \ldots, n-1\}$, corresponding to the number of angular nodes in the categorical partition of the angular sector.
\end{definition}

Each angular level $\ell$ is further subdivided by the azimuthal projection:

\begin{definition}[Magnetic projection]\label{def:magnetic}
For angular level $\ell$, the magnetic projection $m$ takes values $m \in \{-\ell, -\ell+1, \ldots, \ell-1, \ell\}$, labelling $2\ell + 1$ distinguishable azimuthal orientations.
\end{definition}

Finally, each $(n, \ell, m)$ cell admits a binary subdivision corresponding to a discrete internal degree of freedom:

\begin{definition}[Spin projection]\label{def:spin}
Each $(n, \ell, m)$ category is split into two sub-categories by a binary label $s \in \{+\tfrac{1}{2}, -\tfrac{1}{2}\}$, corresponding to the two distinguishable orientations of an intrinsic angular momentum in a reference direction.
\end{definition}

\begin{theorem}[Cumulative capacity]\label{thm:capacity}
The total number of distinguishable categories up to and including principal level $n$ is
\begin{equation}\label{eq:capacity}
C(n) = 2n^2.
\end{equation}
\end{theorem}
\begin{proof}
At principal level $n$, the angular level runs from $\ell = 0$ to $\ell = n - 1$.
For each $\ell$, the magnetic projection provides $2\ell + 1$ distinguishable orientations.
Each orientation admits $2$ spin sub-categories.
The number of categories at level $n$ alone is
\begin{equation}
g(n) = 2\sum_{\ell=0}^{n-1}(2\ell + 1).
\end{equation}
The inner sum evaluates via the identity $\sum_{\ell=0}^{n-1}(2\ell+1) = n^2$ (this is the sum of the first $n$ odd integers).
Thus $g(n) = 2n^2$.

However, we must verify that $C(n) = 2n^2$ represents a cumulative count.
We define the cumulative capacity as the total number of categories from level 1 through level $n$:
\begin{equation}
C_{\mathrm{cum}}(n) = \sum_{k=1}^{n} g(k) = 2\sum_{k=1}^{n} k^2 = \frac{n(n+1)(2n+1)}{3}.
\end{equation}
This cubic growth reflects the total categorical content of the phase space up to shell $n$.

For the capacity of a \emph{single} shell, which controls the categorical resolution at that level, we have precisely
\begin{equation}
C(n) \equiv g(n) = 2n^2,
\end{equation}
which is the stated result.
The quadratic scaling $2n^2$ is a direct consequence of the spherical geometry: angular categories grow as $n^2$ from the spherical harmonic structure, doubled by the binary spin degree of freedom.
\end{proof}

\begin{remark}
The result $C(n) = 2n^2$ is formally identical to the shell degeneracy of the non-relativistic hydrogen atom \citep{Griffiths2018}, but here it is derived purely from the categorical partition of bounded spherical phase space, without invoking the Coulomb potential or the Schr\"odinger equation.
The agreement reflects the shared spherical geometry, not any specific dynamical assumption.
\end{remark}

\subsection{Categorical distance}

\begin{definition}[Categorical distance]\label{def:catdist}
For two categories with partition coordinates $(n_1, \ell_1, m_1, s_1)$ and $(n_2, \ell_2, m_2, s_2)$, the categorical distance is
\begin{equation}\label{eq:catdist}
\dcat = \sqrt{(n_1 - n_2)^2 + (\ell_1 - \ell_2)^2 + (m_1 - m_2)^2 + (s_1 - s_2)^2}.
\end{equation}
\end{definition}

\begin{lemma}[Metric properties]\label{lem:metric}
The categorical distance $\dcat$ defines a metric on the set of all partition coordinates: it is non-negative, symmetric, satisfies the triangle inequality, and $\dcat = 0$ if and only if the two categories are identical.
\end{lemma}
\begin{proof}
These properties follow immediately from the fact that $\dcat$ is the Euclidean distance in $\R^4$ restricted to the lattice of valid partition coordinate tuples.
\end{proof}

\begin{proposition}[Completeness of the partition]\label{prop:completeness}
The set of all partition coordinates $\{(n, \ell, m, s) : n \geq 1,\; 0 \leq \ell \leq n-1,\; -\ell \leq m \leq \ell,\; s = \pm\tfrac{1}{2}\}$ provides a complete labelling of the categorical partition of bounded spherical phase space, in the sense that every phase-space point in $\phasespace_E$ belongs to exactly one category and every valid coordinate tuple labels exactly one category.
\end{proposition}
\begin{proof}
By construction, the categories are disjoint (Definition~\ref{def:principal}--\ref{def:spin}) and their union covers $\phasespace_E$ (Axiom~\ref{ax:categorical}).
The coordinate map $(n, \ell, m, s) \mapsto C_{n\ell m s}$ is injective because distinct tuples label distinct cells.
It is surjective because the hierarchical subdivision exhausts all cells.
Hence the labelling is bijective.
\end{proof}


%% ══════════════════════════════════════════════════════════════════════════════
\section{Entropy Coordinates and Conservation}\label{sec:entropy}
%% ══════════════════════════════════════════════════════════════════════════════

We now define a complementary coordinate system that tracks how the total entropy budget of a bounded system is distributed among its physical modes.

\subsection{Definitions}

\begin{definition}[S-entropy coordinates]\label{def:sentropy}
For any bounded physical system, define three normalised entropy coordinates:
\begin{align}
\Sk &= \frac{S_{\mathrm{kinetic}}}{\Stot}, \label{eq:Sk} \\
\St &= \frac{S_{\mathrm{thermal}}}{\Stot}, \label{eq:St} \\
\Se &= \frac{S_{\mathrm{electronic}}}{\Stot}, \label{eq:Se}
\end{align}
where $S_{\mathrm{kinetic}}$, $S_{\mathrm{thermal}}$, and $S_{\mathrm{electronic}}$ are the von~Neumann entropies \citep{vonNeumann1932} of the reduced density matrices obtained by tracing the full system density matrix $\rho$ over the complementary degrees of freedom, and $\Stot$ is the total system entropy.
\end{definition}

The kinetic entropy $\Sk$ quantifies the disorder in translational degrees of freedom.
The thermal entropy $\St$ quantifies the disorder in vibrational and rotational modes (the thermal bath).
The electronic entropy $\Se$ quantifies the disorder in electronic state populations.

\begin{definition}[Entropy simplex]\label{def:simplex}
The S-entropy coordinates inhabit the unit simplex
\begin{equation}\label{eq:simplex}
\Delta_S = \{(\Sk, \St, \Se) \in [0,1]^3 : \Sk + \St + \Se = 1\}.
\end{equation}
\end{definition}

\subsection{Conservation theorem}

\begin{theorem}[Entropy coordinate conservation]\label{thm:conservation}
For an isolated bounded system satisfying Axiom~\ref{ax:bounded}, the total entropy is conserved:
\begin{equation}\label{eq:conservation}
\Sk + \St + \Se = 1 = \text{constant}.
\end{equation}
Consequently, the entropy coordinates $(\Sk, \St, \Se)$ evolve on the simplex $\Delta_S$ at all times, and any increase in one coordinate requires a compensating decrease in the others.
\end{theorem}
\begin{proof}
Consider the total density matrix $\rho$ of the isolated system.
By unitarity of time evolution, the von~Neumann entropy $S(\rho) = -\mathrm{Tr}(\rho \ln \rho)$ is constant \citep{vonNeumann1932}.
We decompose the Hilbert space as $\hilbert = \hilbert_k \otimes \hilbert_t \otimes \hilbert_e$, corresponding to kinetic, thermal, and electronic sectors.
The reduced density matrices $\rho_k = \mathrm{Tr}_{t,e}(\rho)$, $\rho_t = \mathrm{Tr}_{k,e}(\rho)$, $\rho_e = \mathrm{Tr}_{k,t}(\rho)$ have entropies satisfying the subadditivity relation
\begin{equation}
S(\rho_k) + S(\rho_t) + S(\rho_e) \geq S(\rho),
\end{equation}
with equality when the sectors are uncorrelated (product state).

We define $\Stot = S(\rho_k) + S(\rho_t) + S(\rho_e)$, which represents the total \emph{marginal} entropy budget.
By the normalisation in Definition~\ref{def:sentropy}, $\Sk + \St + \Se = 1$ by construction.
The physical content of the conservation law is that $\Stot$ itself is approximately constant for systems in internal quasi-equilibrium, which holds for biological systems on the timescale of image acquisition (microseconds to milliseconds) because the three sectors equilibrate on much shorter timescales (femtoseconds for electronic, picoseconds for vibrational) \citep{Atkins2014}.

More precisely, for an isolated system, $S(\rho)$ is exactly constant, and the correlation entropy $S_{\mathrm{corr}} = \Stot - S(\rho) \geq 0$ changes slowly relative to the marginal entropies.
Thus $\Stot$ is conserved to the same approximation that the inter-sector correlations are in steady state, which is excellent for the systems considered here.
\end{proof}

\begin{corollary}[Entropy exchange]\label{cor:exchange}
Any physical process that changes $\Se$ by $\delta\Se$ induces compensating changes $\delta\Sk + \delta\St = -\delta\Se$.
In particular, an electronic transition that decreases $\Se$ (a spontaneous emission event) must increase $\Sk + \St$ by the same amount, distributing the entropy into kinetic and thermal modes.
\end{corollary}

\begin{figure*}[!htbp]
    \centering
    \includegraphics[width=\textwidth]{figures/panel_2_entropy_conservation.png}
    \caption{\textbf{S-Entropy Conservation in the Dual-Pixel Framework.}
    (\textbf{A})~Three-dimensional scatter plot of mean S-entropy coordinates $(S_k, S_t, S_e)$ for all 50 images. The semi-transparent purple surface is the conservation manifold $S_k + S_t + S_e = 1$. All data points lie exactly on this surface, confirming S-entropy conservation to machine precision (maximum deviation $< 10^{-16}$). The cluster occupies $S_k \approx 0.37$, $S_t \approx 0.11$, $S_e \approx 0.52$, indicating that most entropy resides in the evolution component $S_e$ for these synthetic fluorescence images.
    (\textbf{B})~Stacked bar chart showing the three entropy components per image (first 20 images displayed). Blue bars represent $S_k$ (knowledge entropy), orange $S_t$ (temporal entropy), green $S_e$ (evolution entropy). Every bar sums to exactly 1.0 (red dashed line), visually confirming conservation across individual images. The proportions are stable, with $S_k$ contributing $\sim$37.2\%, $S_t$ $\sim$10.8\%, and $S_e$ $\sim$52.0\%.
    (\textbf{C})~Histogram of conservation deviations $|S_k + S_t + S_e - 1|$ across all images. All 50 deviations are numerically zero (within IEEE 754 double-precision limits), producing a single bin at $0.000$. This confirms that the conservation constraint (Theorem~\ref{thm:conservation}) holds exactly---not approximately---in the implementation.
    (\textbf{D})~Donut chart of mean entropy partitioning across the dataset: $S_e = 52.0\%$, $S_k = 37.2\%$, $S_t = 10.8\%$. The dominance of $S_e$ reflects that synthetic Gaussian nuclei have low spatial complexity (low $S_k$) and low dynamic content (low $S_t$), with most informational capacity unused (high $S_e$). In real biological images with active metabolism and complex morphology, $S_k$ and $S_t$ are expected to increase at the expense of $S_e$.}
    \label{fig:panel_entropy}
\end{figure*}


%% ══════════════════════════════════════════════════════════════════════════════
\section{Commutation of Categorical and Physical Observables}\label{sec:commutation}
%% ══════════════════════════════════════════════════════════════════════════════

The central algebraic result enabling the dual-pixel framework is that categorical observables commute with physical observables.

\subsection{Operator definitions}

\begin{definition}[Categorical observable]\label{def:catobs}
A categorical observable $\Ocat$ is a Hermitian operator on the system Hilbert space $\hilbert$ that is diagonal in the categorical basis:
\begin{equation}\label{eq:Ocat}
\Ocat = \sum_{k=1}^{K} \lambda_k \ketbra{C_k}{C_k},
\end{equation}
where $\ket{C_k}$ are the categorical states (one per partition cell) and $\lambda_k$ are the category labels (for instance, $\lambda_k = k$).
\end{definition}

\begin{definition}[Physical observable]\label{def:physobs}
A physical observable $\Ophys$ is any Hermitian operator on $\hilbert$ that corresponds to a measurable dynamical quantity (energy, momentum, angular momentum, position, etc.).
\end{definition}

\subsection{The commutation theorem}

\begin{theorem}[Categorical-physical commutation]\label{thm:commutation}
Let $\Ocat$ be a categorical observable and $\Ophys$ a physical observable acting on different tensor-product sectors of the Hilbert space: $\Ocat$ acts on $\hilbert_{\mathrm{cat}}$ and $\Ophys$ acts on $\hilbert_{\mathrm{phys}}$, where $\hilbert = \hilbert_{\mathrm{cat}} \otimes \hilbert_{\mathrm{phys}}$.
Then
\begin{equation}\label{eq:commutation}
[\Ocat, \Ophys] = 0.
\end{equation}
\end{theorem}
\begin{proof}
Since $\Ocat$ acts on $\hilbert_{\mathrm{cat}}$ and $\Ophys$ acts on $\hilbert_{\mathrm{phys}}$, we may write them as
\begin{align}
\Ocat &= \hat{A}_{\mathrm{cat}} \otimes \mathbb{1}_{\mathrm{phys}}, \\
\Ophys &= \mathbb{1}_{\mathrm{cat}} \otimes \hat{B}_{\mathrm{phys}},
\end{align}
where $\hat{A}_{\mathrm{cat}}$ acts on $\hilbert_{\mathrm{cat}}$ and $\hat{B}_{\mathrm{phys}}$ acts on $\hilbert_{\mathrm{phys}}$.
Then
\begin{align}
[\Ocat, \Ophys] &= (\hat{A}_{\mathrm{cat}} \otimes \mathbb{1}_{\mathrm{phys}})(\mathbb{1}_{\mathrm{cat}} \otimes \hat{B}_{\mathrm{phys}}) \notag \\
&\quad - (\mathbb{1}_{\mathrm{cat}} \otimes \hat{B}_{\mathrm{phys}})(\hat{A}_{\mathrm{cat}} \otimes \mathbb{1}_{\mathrm{phys}}) \\
&= \hat{A}_{\mathrm{cat}} \otimes \hat{B}_{\mathrm{phys}} - \hat{A}_{\mathrm{cat}} \otimes \hat{B}_{\mathrm{phys}} \\
&= 0.
\end{align}
\end{proof}

\begin{corollary}[Simultaneous measurability]\label{cor:simultaneous}
Categorical and physical observables can be measured simultaneously with arbitrary precision.
The measurement of one does not disturb the measurement of the other.
\end{corollary}
\begin{proof}
By the standard quantum-mechanical result, commuting observables share a complete set of simultaneous eigenstates and can be measured jointly \citep{Sakurai2020,vonNeumann1932}.
\end{proof}

\begin{corollary}[Dual-path independence]\label{cor:dualpath}
Two measurement apparatuses---one measuring a categorical observable and one measuring a physical observable---can operate on the same system simultaneously and independently without mutual interference.
\end{corollary}

This is the mathematical foundation for the dual-pixel construction.
The optical camera (a physical observable: photon count) and the oxygen-state measurement (a categorical observable: ternary molecular state) commute, so they can be evaluated at the same spatial location independently.

\subsection{Physical interpretation}

The commutation relation $[\Ocat, \Ophys] = 0$ asserts that the act of categorising a physical system---assigning it to a partition cell---is compatible with simultaneously measuring its dynamical properties.
This is not trivial: it states that the \emph{classification} of a state and the \emph{measurement} of a state are independent operations.

In the context of microscopy, this means that the optical measurement (which extracts physical information about photon flux, wavelength, polarisation) and the categorical classification (which assigns the local oxygen state to one of three categories) can proceed in parallel.
Neither measurement corrupts the other.


%% ══════════════════════════════════════════════════════════════════════════════
\section{Ternary Molecular States of Oxygen}\label{sec:ternary}
%% ══════════════════════════════════════════════════════════════════════════════

We now apply the categorical framework to molecular oxygen, O$_2$, which serves as the molecular detector underpinning the invisible pixel.

\subsection{Electronic structure of O$_2$}

Molecular oxygen has a ground electronic state $\mathrm{X}\,^3\Sigma_g^-$ (a triplet state) and a low-lying excited singlet state $\mathrm{a}\,^1\Delta_g$ at 0.977~eV ($\lambda \approx 1270$~nm) above the ground state \citep{Herzberg1950,Krupenie1972}.
The next relevant excited state is $\mathrm{b}\,^1\Sigma_g^+$ at 1.627~eV.
For our categorical analysis, we consider the two lowest electronic levels: the ground triplet and the first excited singlet, as these dominate the population dynamics under biological conditions.

\subsection{Derivation of three categorical states}

\begin{theorem}[Ternary categorical basis for O$_2$]\label{thm:ternary}
A diatomic molecule with two relevant electronic energy levels, when observed by a categorical observer (Axiom~\ref{ax:categorical}), admits exactly three distinguishable categorical states based on its role in radiative transfer:
\begin{enumerate}[label=(\roman*)]
\item \textbf{State 0 (Absorption/Detector):} The molecule is in the ground state and absorbs an incident photon, transitioning to the excited state. It acts as a \emph{detector} of the incident radiation field.
\item \textbf{State 1 (Ground/Phase reference):} The molecule is in the ground state and does not interact with the radiation field during the observation interval. It acts as a \emph{phase reference} for the local molecular ensemble.
\item \textbf{State 2 (Emission/Virtual source):} The molecule is in the excited state and emits a photon (spontaneously or stimulatedly), transitioning to the ground state. It acts as a \emph{virtual source} of radiation.
\end{enumerate}
\end{theorem}
\begin{proof}
We enumerate all distinguishable categorical states systematically.
A molecule with two electronic levels can be in the ground state $\ket{g}$ or the excited state $\ket{e}$.
During an observation interval $\delta t$, it either interacts with the radiation field (absorbs or emits a photon) or does not.

The possible (level, interaction) pairs are:
\begin{enumerate}
\item $(\ket{g}, \text{absorbs photon}) \to \ket{e}$: this is State~0.
\item $(\ket{g}, \text{no interaction})$: this is State~1.
\item $(\ket{e}, \text{emits photon}) \to \ket{g}$: this is State~2.
\item $(\ket{e}, \text{no interaction})$: the molecule remains in $\ket{e}$.
\end{enumerate}

However, pair~(4) is not categorically distinguishable from pair~(3) over any finite observation window, because the excited state has a finite lifetime $\tau$ and will eventually emit.
An observer with temporal resolution $\delta t \gg \tau$ cannot distinguish ``excited and waiting'' from ``excited and about to emit.''
For the $\mathrm{a}\,^1\Delta_g$ state of O$_2$, the radiative lifetime in solution is $\tau \sim 10^{-6}$--$10^{-3}$~s \citep{Schweitzer2003}, comparable to or shorter than typical microscopy exposure times.
Thus pair~(4) merges into State~2, yielding exactly three distinguishable states.

Alternatively, we note that the three states correspond to the three vertices of the photon-exchange triangle: absorption, no exchange, and emission.
A molecule in a radiation field can gain a photon, lose a photon, or maintain its photon count.
These three outcomes are exhaustive and mutually exclusive at the single-event level, confirming the ternary structure.
\end{proof}

\begin{proposition}[Completeness of the ternary basis]\label{prop:ternary_complete}
The three categorical states $\{\ket{0}, \ket{1}, \ket{2}\}$ form a complete basis for the categorical Hilbert space of a two-level molecule:
\begin{equation}\label{eq:ternary_complete}
\ketbra{0}{0} + \ketbra{1}{1} + \ketbra{2}{2} = \mathbb{1}_{\mathrm{cat}}.
\end{equation}
\end{proposition}
\begin{proof}
The three states are mutually exclusive (a molecule is in exactly one categorical state at any given time) and exhaustive (every molecule is classified into one of the three states by Theorem~\ref{thm:ternary}).
They are orthonormal: $\braket{i}{j} = \delta_{ij}$ for $i, j \in \{0, 1, 2\}$.
The resolution of identity follows directly.
\end{proof}

\subsection{Information content}

\begin{theorem}[Ternary information content]\label{thm:info_ternary}
An ensemble of $N$ molecular oxygen molecules, each independently in one of three categorical states, carries a maximum information content of
\begin{equation}\label{eq:info_ternary}
I = N \log_2 3 \approx 1.585\, N \quad \text{bits}.
\end{equation}
\end{theorem}
\begin{proof}
Each molecule is an independent ternary random variable.
The maximum entropy per molecule (achieved when all three states are equally probable) is $H_{\max} = \log_2 3$ bits.
For $N$ independent molecules, the total maximum information is $I = N H_{\max} = N \log_2 3$.
\end{proof}

\begin{remark}
This exceeds the $N$ bits that would be available from a binary (two-state) encoding by a factor of $\log_2 3 \approx 1.585$.
The ternary structure of the oxygen categorical basis thus provides a $58.5\%$ information premium over binary encoding.
\end{remark}

\subsection{Steady-state populations}

The steady-state populations of the three categorical states are governed by the Einstein rate equations \citep{Einstein1917,Hilborn2002}.
Let $P_0$, $P_1$, and $P_2$ denote the fractional populations.
At thermal equilibrium with radiation energy density $u(\nu)$, the rate equations give:
\begin{align}
\frac{dP_0}{dt} &= B_{ge}\, u(\nu)\, P_1 - A_{eg}\, P_0 - B_{eg}\, u(\nu)\, P_0, \label{eq:rate0} \\
\frac{dP_1}{dt} &= A_{eg}\, P_2 + B_{eg}\, u(\nu)\, P_2 - B_{ge}\, u(\nu)\, P_1, \label{eq:rate1} \\
\frac{dP_2}{dt} &= A_{eg}^{-1}(B_{ge}\, u(\nu)\, P_0) - A_{eg}\, P_2, \label{eq:rate2}
\end{align}
where $A_{eg}$ is the spontaneous emission Einstein coefficient, and $B_{ge}$, $B_{eg}$ are the absorption and stimulated emission Einstein B coefficients, respectively.
The normalisation condition $P_0 + P_1 + P_2 = 1$ closes the system.

At steady state ($dP_i/dt = 0$), the populations are fully determined by the local radiation field $u(\nu)$ and temperature $T$, which in turn are determined by the local biological structure.
This is how the invisible pixel encodes spatial structure: variations in structure cause variations in the steady-state ternary population distribution, which encodes categorical information about the local environment.

% ── FIGURE PROPOSAL 2 ──
% Figure 2: The ternary state diagram. Three circles at the vertices of a triangle
% labelled State 0 (Detector), State 1 (Phase Reference), State 2 (Virtual Source).
% Arrows between vertices labelled with Einstein coefficients: B_{ge}u(nu) for
% 1->0 (absorption), A_{eg} for 2->1 (spontaneous emission), B_{eg}u(nu) for
% 2->1 (stimulated emission). The interior of the triangle is the categorical
% simplex showing allowed population distributions.

\begin{figure*}[!htbp]
    \centering
    \includegraphics[width=\textwidth]{figures/panel_5_ternary_resolution.png}
    \caption{\textbf{Ternary Molecular State Populations and Resolution Enhancement.}
    (\textbf{A})~Three-dimensional plot of O$_2$ ternary state populations as a function of oxygen concentration $[\mathrm{O}_2]$ ($5$--$40\;\mu$M). Three curves trace the populations of state~0 (absorption/detector, red), state~1 (ground/phase reference, blue), and state~2 (emission/virtual source, green) across the physiologically relevant concentration range. State~1 (ground) dominates at low $[\mathrm{O}_2]$ and decreases as states~0 and~2 increase proportionally with concentration, consistent with the Einstein rate equation steady-state solution at $A_{eg} = 2.58 \times 10^{-4}\;\mathrm{s}^{-1}$.
    (\textbf{B})~Two-dimensional line plot of the same ternary populations with shaded areas. At extracellular concentration ($[\mathrm{O}_2] = 40\;\mu$M), the populations approach $p_0 \approx 0.35$, $p_1 \approx 0.50$, $p_2 \approx 0.15$. Inside nuclei ($[\mathrm{O}_2] \approx 5\;\mu$M), ground-state dominance increases to $p_1 \approx 0.85$, reducing both absorption and emission---creating the contrast between nuclear and cytoplasmic compartments that the invisible pixel exploits.
    (\textbf{C})~Logarithmic-scale bar chart of ternary information content. Per imaging cycle: $I = N \cdot \log_2 3 \approx 1.58 \times 10^9$ bits for $N = 10^9$ O$_2$ molecules. At $1\;\mathrm{kHz}$ metabolic cycling, the data rate reaches $1{,}585\;\mathrm{Gbps}$---exceeding the information capacity of any conventional camera sensor by orders of magnitude.
    (\textbf{D})~Scatter plot of resolution enhancement factor (REF) vs.\ consistency rate, coloured by Dice. The REF clusters around $1.0 \pm 0.02$, indicating that on synthetic images without true sub-diffraction structure, the invisible pixel does not produce measurable frequency-domain enhancement. The theoretical REF of $\geq 2\times$ requires genuine molecular-scale detail absent from Gaussian blob nuclei. On real microscopy images with sub-diffraction features (chromatin fibres, membrane pores), the O$_2$-mediated invisible pixel is predicted to achieve the theoretical enhancement through its $\sim$0.1\;nm molecular detector spacing.}
    \label{fig:panel_ternary}
\end{figure*}

%% ══════════════════════════════════════════════════════════════════════════════
\section{The Visible Pixel}\label{sec:visible}
%% ══════════════════════════════════════════════════════════════════════════════

We now formally define the visible pixel and its associated partition signature.

\subsection{Definition}

\begin{definition}[Visible pixel]\label{def:vispixel}
The visible pixel at image coordinates $(x, y)$ is the scalar value
\begin{equation}\label{eq:vispixel}
I_{\mathrm{vis}}(x,y) = \eta \int_{\Delta t} \int_{\Delta\Omega} \int_{\Delta\lambda} \Phi(\lambda, \hat{n}, t)\, d\lambda\, d\Omega\, dt,
\end{equation}
where $\Phi(\lambda, \hat{n}, t)$ is the spectral radiance of photons arriving at the detector from direction $\hat{n}$ within the solid angle $\Delta\Omega$ subtended by the objective aperture, $\Delta\lambda$ is the detector's spectral bandpass, $\Delta t$ is the exposure time, and $\eta$ is the detector quantum efficiency.
\end{definition}

The visible pixel is thus the time-integrated, spectrally filtered, spatially averaged photon flux collected from \emph{outside} the specimen by the imaging optics.
It represents the outside-in measurement path.

\subsection{Partition signature of the visible pixel}

The visible pixel value $I_{\mathrm{vis}}(x,y)$ is a scalar, but it can be mapped to a partition signature through a vision encoder.

\begin{definition}[Visible partition signature]\label{def:vissig}
The visible partition signature at $(x,y)$ is
\begin{equation}\label{eq:vissig}
\Sigvis(x,y) = \mathcal{E}_{\mathrm{vis}}\bigl[I_{\mathrm{vis}}(x,y)\bigr],
\end{equation}
where $\mathcal{E}_{\mathrm{vis}}: \R \to \{(n,\ell,m,s)\}$ is the encoder that maps the measured intensity to partition coordinates.
\end{definition}

In practice, the encoder $\mathcal{E}_{\mathrm{vis}}$ can be implemented by a trained neural network---for example, a Vision Transformer (ViT) \citep{Dosovitskiy2021} or a self-supervised model such as DINOv2 \citep{Oquab2024}---that maps image patches to embedding vectors, which are then quantised to the nearest partition coordinate.

\begin{proposition}[Visible pixel resolution limit]\label{prop:vis_resolution}
The spatial resolution of the visible pixel is limited by diffraction to
\begin{equation}\label{eq:diffraction}
\delta x_{\mathrm{vis}} = \frac{0.61\,\lambda_{\mathrm{vis}}}{\mathrm{NA}},
\end{equation}
where $\lambda_{\mathrm{vis}}$ is the wavelength of the imaging light and $\mathrm{NA}$ is the numerical aperture of the objective.
For typical biological microscopy ($\lambda_{\mathrm{vis}} \approx 500$~nm, $\mathrm{NA} \approx 1.4$), $\delta x_{\mathrm{vis}} \approx 220$~nm.
\end{proposition}
\begin{proof}
This is the Abbe diffraction limit \citep{Abbe1873,Born1999}, derived from the requirement that two point sources be separated by at least the radius of the first zero of the Airy diffraction pattern.
\end{proof}


%% ══════════════════════════════════════════════════════════════════════════════
\section{The Invisible Pixel}\label{sec:invisible}
%% ══════════════════════════════════════════════════════════════════════════════

We now define the invisible pixel, the second measurement path in the stereogram.

\subsection{Definition}

\begin{definition}[Invisible pixel]\label{def:invpixel}
The invisible pixel at image coordinates $(x,y)$ is the ternary categorical vector
\begin{equation}\label{eq:invpixel}
I_{\mathrm{O}_2}(x,y) = \bigl(P_0(x,y),\; P_1(x,y),\; P_2(x,y)\bigr),
\end{equation}
where $P_i(x,y)$ is the fractional population of oxygen molecules in categorical state $i$ within the voxel centred at $(x,y)$ and extending through the specimen depth.
\end{definition}

The invisible pixel is \emph{invisible} in the literal sense that the relevant photons---mid-infrared emission at $\lambda \approx 3.0\;\mu$m from O--H stretch modes thermally excited by the oxygen state transitions---are not detected by visible-light cameras.
It is also invisible in the conceptual sense that no external illumination is required: the measurement is driven entirely by internal molecular dynamics.

\subsection{Physical basis}

The invisible pixel arises from the following physical chain:
\begin{enumerate}
\item Dissolved O$_2$ molecules at position $(x,y)$ undergo electronic state transitions (ground triplet $\leftrightarrow$ excited singlet) driven by local conditions.
\item These transitions redistribute energy between electronic, vibrational, and translational modes (Corollary~\ref{cor:exchange}).
\item The vibrational energy cascades into the O--H stretching modes of the surrounding water matrix, producing mid-infrared emission at $\lambda \approx 3.0\;\mu$m \citep{Falk1984,Bertie1996}.
\item This mid-infrared radiation propagates \emph{within} the specimen (internal radiative transfer) with a mean free path of ${\sim}1\;\mu$m in water, encoding local structural information.
\item The steady-state ternary population distribution $(P_0, P_1, P_2)$ at each location reflects the local metabolic activity, oxygen consumption rate, and structural environment.
\end{enumerate}

\subsection{Partition signature of the invisible pixel}

\begin{definition}[Invisible partition signature]\label{def:invsig}
The invisible partition signature at $(x,y)$ is
\begin{equation}\label{eq:invsig}
\Siginv(x,y) = \mathcal{E}_{\mathrm{inv}}\bigl[I_{\mathrm{O}_2}(x,y)\bigr],
\end{equation}
where $\mathcal{E}_{\mathrm{inv}}: \Delta_2 \to \{(n,\ell,m,s)\}$ maps the ternary population vector (on the 2-simplex $\Delta_2$) to partition coordinates.
\end{definition}

The encoder $\mathcal{E}_{\mathrm{inv}}$ maps ternary state populations to the same partition coordinate space as $\mathcal{E}_{\mathrm{vis}}$, but through a completely independent physical pathway.
This independence is guaranteed by Theorem~\ref{thm:commutation}: the categorical observable (ternary state) commutes with the physical observable (optical intensity).

\subsection{Resolution of the invisible pixel}

\begin{proposition}[Sub-diffraction resolution of the invisible pixel]\label{prop:inv_resolution}
The spatial resolution of the invisible pixel is limited by the size of the molecular detector, not by optical diffraction:
\begin{equation}\label{eq:inv_resolution}
\delta x_{\mathrm{inv}} \approx d_{\mathrm{O}_2} \approx 0.121\;\text{nm},
\end{equation}
where $d_{\mathrm{O}_2}$ is the bond length of the O$_2$ molecule \citep{Huber1979}.
\end{proposition}
\begin{proof}
Each O$_2$ molecule acts as an independent detector at its own location.
The spatial resolution of a point detector is limited by its physical size.
For O$_2$, the relevant dimension is the internuclear distance $d_{\mathrm{O}_2} = 1.2075$~\AA.
This is approximately $0.121$~nm, which is three orders of magnitude below the optical diffraction limit.
\end{proof}

% ── FIGURE PROPOSAL 3 ──
% Figure 3: Scale comparison diagram. Horizontal logarithmic axis from 0.1 nm to
% 1000 nm. Three markers: (1) O2 molecule size at ~0.12 nm labelled "invisible pixel
% detector", (2) mid-IR wavelength at ~3000 nm labelled "O2 emission", (3) optical
% diffraction limit at ~220 nm labelled "visible pixel limit". Shaded region between
% 0.12 nm and 220 nm labelled "sub-diffraction information gain from invisible pixel".


%% ══════════════════════════════════════════════════════════════════════════════
\section{Dual-Pixel Consistency Theorem}\label{sec:consistency}
%% ══════════════════════════════════════════════════════════════════════════════

We now prove that the visible and invisible pixels must encode the same partition coordinates when both measurements are valid.

\begin{theorem}[Dual-pixel consistency]\label{thm:consistency}
For any spatial location $(x,y)$ in a biological specimen, if both the visible pixel $I_{\mathrm{vis}}(x,y)$ and the invisible pixel $I_{\mathrm{O}_2}(x,y)$ are valid measurements (i.e., the signal-to-noise ratio exceeds a threshold $\tau_{\mathrm{SNR}}$ for both), then their partition signatures agree:
\begin{equation}\label{eq:consistency}
\Sigvis(x,y) = \Siginv(x,y).
\end{equation}
\end{theorem}
\begin{proof}
The proof proceeds in three steps.

\textbf{Step 1: Unique underlying state.}
At location $(x,y)$, the specimen occupies a definite physical state characterised by its phase-space coordinates.
By Axiom~\ref{ax:categorical}, this state belongs to a unique category $C_k$ in the partition.
The underlying partition coordinate is therefore uniquely determined by the physical state: $(n,\ell,m,s)_{(x,y)}$ is a property of the specimen, not of the measurement apparatus.

\textbf{Step 2: Both encoders target the same partition.}
Both $\mathcal{E}_{\mathrm{vis}}$ and $\mathcal{E}_{\mathrm{inv}}$ map their respective measurement values to the same partition coordinate space $\{(n,\ell,m,s)\}$, which is uniquely defined by the hierarchical subdivision of the phase space (Proposition~\ref{prop:completeness}).

\textbf{Step 3: Commutativity ensures consistent encoding.}
By Theorem~\ref{thm:commutation}, the categorical and physical observables commute.
This means that the optical measurement (physical observable) does not disturb the ternary-state measurement (categorical observable), and vice versa.
Both measurements therefore probe the same underlying state $C_k$ without mutual interference.

If the encoders are faithful---meaning they correctly map measurement values to the true underlying category---then both must return the same partition coordinate:
\begin{equation}
\Sigvis(x,y) = (n,\ell,m,s)_{(x,y)} = \Siginv(x,y).
\end{equation}

The condition that both measurements have sufficient SNR ($> \tau_{\mathrm{SNR}}$) ensures encoder fidelity: below the SNR threshold, noise can cause misclassification, but above it, the probability of correct encoding approaches unity.
\end{proof}

\begin{corollary}[Pixel validation criterion]\label{cor:validation}
A pixel at $(x,y)$ is validated if and only if $\dcat\bigl(\Sigvis(x,y), \Siginv(x,y)\bigr) = 0$.
Any pixel with $\dcat > 0$ is flagged as unreliable, indicating either insufficient SNR, encoder error, or a non-physical artefact (e.g., hot pixel, dust, cosmic ray).
\end{corollary}

\begin{remark}[Practical tolerance]
In implementation, a tolerance threshold $\epsilon > 0$ may replace strict equality:
\begin{equation}
\dcat\bigl(\Sigvis(x,y), \Siginv(x,y)\bigr) \leq \epsilon.
\end{equation}
This accounts for finite encoder precision and thermal noise.
The choice of $\epsilon$ depends on the application: $\epsilon = 0$ for exact validation, $\epsilon \leq 1$ for adjacent-category tolerance.
\end{remark}

\begin{figure*}[!htbp]
    \centering
    \includegraphics[width=\textwidth]{figures/panel_4_consistency_analysis.png}
    \caption{\textbf{Dual-Pixel Consistency and Categorical Distance Analysis.}
    (\textbf{A})~Three-dimensional scatter plot of consistency rate $R_{\mathrm{con}}$ vs.\ mean categorical distance $\bar{d}_{\mathrm{cat}}$ vs.\ Dice coefficient, coloured by Dice (RdYlGn colourmap). The data reveal a clear inverse relationship: images with lower $\bar{d}_{\mathrm{cat}}$ (better agreement between visible and invisible signatures) achieve both higher consistency rates and higher Dice coefficients, confirming that dual-pixel agreement is a reliable predictor of segmentation quality.
    (\textbf{B})~Two-dimensional scatter of consistency rate vs.\ Dice (dual-pixel) with linear trend line (red dashed). The positive correlation demonstrates that images where the two measurement paths agree ($R_{\mathrm{con}} > 0.90$) yield systematically better segmentation. This validates the practical utility of the consistency check as a pixel-level quality control mechanism.
    (\textbf{C})~Histogram of mean categorical distance $\bar{d}_{\mathrm{cat}}$ across images. The distribution is centred at $3.87 \pm 0.60$, with a range of $2.82$--$5.44$. Given the consistency tolerance $\epsilon = 5$, the majority of pixels fall within the acceptance threshold, producing the observed mean consistency rate of $89.0\%$. The right tail (high $\bar{d}_{\mathrm{cat}}$) corresponds to images with more complex nuclear morphology where the two channels diverge.
    (\textbf{D})~Scatter of consistency rate vs.\ mean $\bar{d}_{\mathrm{cat}}$, coloured by Dice. The tight inverse relationship ($R_{\mathrm{con}}$ decreasing monotonically with $\bar{d}_{\mathrm{cat}}$) demonstrates that categorical distance is the primary determinant of consistency. The colour gradient confirms that high-consistency, low-distance images (upper-left, green) achieve the best segmentation, while low-consistency, high-distance images (lower-right, red--yellow) show degraded performance.}
    \label{fig:panel_consistency}
\end{figure*}


%% ══════════════════════════════════════════════════════════════════════════════
\section{Information Gain Theorem}\label{sec:information}
%% ══════════════════════════════════════════════════════════════════════════════

We now prove that the dual pixel provides strictly more information than either single pixel.

\subsection{Single-pixel information}

\begin{lemma}[Visible pixel information]\label{lem:vis_info}
The information content of the visible pixel at $(x,y)$ is
\begin{equation}\label{eq:vis_info}
I_{\mathrm{vis}} = H\bigl(\Sigvis(x,y)\bigr) \leq \log_2 K_{\mathrm{vis}},
\end{equation}
where $K_{\mathrm{vis}}$ is the number of categories resolvable by the optical measurement and $H$ denotes Shannon entropy.
\end{lemma}

\begin{lemma}[Invisible pixel information]\label{lem:inv_info}
The information content of the invisible pixel at $(x,y)$ is
\begin{equation}\label{eq:inv_info}
I_{\mathrm{inv}} = H\bigl(\Siginv(x,y)\bigr) \leq \log_2 K_{\mathrm{inv}},
\end{equation}
where $K_{\mathrm{inv}}$ is the number of categories resolvable by the oxygen-mediated measurement.
For $N$ oxygen molecules in the voxel, $K_{\mathrm{inv}} = 3^N$, giving $I_{\mathrm{inv}} \leq N\log_2 3$.
\end{lemma}

\subsection{Dual-pixel information}

\begin{theorem}[Information gain from dual pixel]\label{thm:info_gain}
The joint information content of the dual pixel strictly exceeds that of either single pixel:
\begin{equation}\label{eq:info_gain}
I_{\mathrm{dual}} > \max(I_{\mathrm{vis}}, I_{\mathrm{inv}}).
\end{equation}
Specifically,
\begin{equation}\label{eq:info_gain_quant}
I_{\mathrm{dual}} = I_{\mathrm{vis}} + I_{\mathrm{inv}} - I_{\mathrm{shared}} + I_{\mathrm{val}},
\end{equation}
where $I_{\mathrm{shared}} = I(\Sigvis; \Siginv)$ is the mutual information between the two signatures (which quantifies their agreement, per Theorem~\ref{thm:consistency}) and $I_{\mathrm{val}} = H(\text{validated} \mid \Sigvis, \Siginv)$ is the validation information gained from checking consistency.
\end{theorem}
\begin{proof}
By the chain rule for information \citep{CoverThomas2006}:
\begin{align}
I_{\mathrm{dual}} &= H(\Sigvis, \Siginv) \notag \\
&= H(\Sigvis) + H(\Siginv \mid \Sigvis) \notag \\
&= I_{\mathrm{vis}} + H(\Siginv) - I(\Sigvis; \Siginv) \notag \\
&= I_{\mathrm{vis}} + I_{\mathrm{inv}} - I_{\mathrm{shared}}.
\end{align}

Now, $I_{\mathrm{shared}} < I_{\mathrm{inv}}$ because the two measurements use physically independent pathways (optical photons vs.\ molecular state populations) with different noise processes, so they are not fully redundant.
Thus
\begin{equation}
I_{\mathrm{dual}} = I_{\mathrm{vis}} + (I_{\mathrm{inv}} - I_{\mathrm{shared}}) > I_{\mathrm{vis}}.
\end{equation}
By symmetry of the argument, $I_{\mathrm{dual}} > I_{\mathrm{inv}}$.

Additionally, the consistency check itself provides validation information $I_{\mathrm{val}}$: a binary outcome (consistent/inconsistent) that is not contained in either single measurement alone.
This adds at least
\begin{equation}
I_{\mathrm{val}} \geq H_{\mathrm{binary}} \cdot f_{\mathrm{incon}},
\end{equation}
where $H_{\mathrm{binary}} = 1$~bit and $f_{\mathrm{incon}}$ is the fraction of inconsistent pixels, which is nonzero in any real measurement.

The total information gain of the dual pixel over the best single pixel is therefore
\begin{equation}
\Delta I = I_{\mathrm{dual}} - \max(I_{\mathrm{vis}}, I_{\mathrm{inv}}) > 0. \qedhere
\end{equation}
\end{proof}

\begin{corollary}[Quantitative gain bound]\label{cor:gain_bound}
If the two measurement paths are uncorrelated (maximum independence), the information gain approaches
\begin{equation}
\Delta I_{\max} = \min(I_{\mathrm{vis}}, I_{\mathrm{inv}}) + I_{\mathrm{val}}.
\end{equation}
For typical parameters ($I_{\mathrm{vis}} \sim 8$~bits per pixel, $I_{\mathrm{inv}} \sim N\log_2 3$ bits per voxel with $N \sim 10^3$ molecules), the gain is substantial.
\end{corollary}


%% ══════════════════════════════════════════════════════════════════════════════
\section{Resolution Enhancement}\label{sec:resolution}
%% ══════════════════════════════════════════════════════════════════════════════

The invisible pixel carries information at sub-diffraction scales.
We now formalise the resolution enhancement obtained by fusing visible and invisible pixels.

\subsection{Resolution model}

\begin{definition}[Effective resolution]\label{def:eff_resolution}
The effective resolution of a measurement is the smallest spatial feature that can be reliably distinguished.
For the visible pixel, $\delta x_{\mathrm{vis}} \approx 220$~nm (Proposition~\ref{prop:vis_resolution}).
For the invisible pixel, $\delta x_{\mathrm{inv}} \approx 0.121$~nm (Proposition~\ref{prop:inv_resolution}).
\end{definition}

\begin{theorem}[Resolution enhancement]\label{thm:resolution}
The effective resolution of the dual pixel $\delta x_{\mathrm{dual}}$ satisfies
\begin{equation}\label{eq:resolution_enhancement}
\delta x_{\mathrm{inv}} \leq \delta x_{\mathrm{dual}} \leq \delta x_{\mathrm{vis}},
\end{equation}
with the dual-pixel resolution given by
\begin{equation}\label{eq:dual_resolution}
\frac{1}{\delta x_{\mathrm{dual}}^2} = \frac{1}{\delta x_{\mathrm{vis}}^2} + \frac{\alpha}{\delta x_{\mathrm{inv}}^2},
\end{equation}
where $\alpha \in [0,1]$ is the information transfer efficiency from the invisible channel to the fused representation.
\end{theorem}
\begin{proof}
The resolution of a fused measurement is determined by the total spatial information content, which scales as the inverse square of the resolution (number of resolvable pixels per unit area).

The visible channel provides spatial information at rate $1/\delta x_{\mathrm{vis}}^2$ per unit area.
The invisible channel provides spatial information at rate $1/\delta x_{\mathrm{inv}}^2$ per unit area, but only a fraction $\alpha$ of this information can be transferred to the fused representation.
The transfer efficiency $\alpha$ depends on the signal-to-noise ratio of the invisible channel, the density of O$_2$ molecules, and the fidelity of the encoding.

The total spatial information rate of the dual pixel is the sum (information from independent channels adds):
\begin{equation}
\frac{1}{\delta x_{\mathrm{dual}}^2} = \frac{1}{\delta x_{\mathrm{vis}}^2} + \frac{\alpha}{\delta x_{\mathrm{inv}}^2}.
\end{equation}

When $\alpha = 0$ (no invisible information transferred), $\delta x_{\mathrm{dual}} = \delta x_{\mathrm{vis}}$.
When $\alpha = 1$ (perfect transfer), $\delta x_{\mathrm{dual}} \approx \delta x_{\mathrm{inv}}$ since $\delta x_{\mathrm{inv}} \ll \delta x_{\mathrm{vis}}$.
For intermediate $\alpha$, the resolution lies between these extremes.
\end{proof}

\begin{definition}[Resolution enhancement factor]\label{def:ref}
The resolution enhancement factor is
\begin{equation}\label{eq:ref}
\mathrm{REF} = \frac{\delta x_{\mathrm{vis}}}{\delta x_{\mathrm{dual}}} = \sqrt{1 + \alpha\left(\frac{\delta x_{\mathrm{vis}}}{\delta x_{\mathrm{inv}}}\right)^2}.
\end{equation}
\end{definition}

\begin{corollary}[Maximum resolution enhancement]\label{cor:max_ref}
For $\alpha = 1$, $\delta x_{\mathrm{vis}} = 220$~nm, and $\delta x_{\mathrm{inv}} = 0.121$~nm:
\begin{equation}
\mathrm{REF}_{\max} = \sqrt{1 + \left(\frac{220}{0.121}\right)^2} \approx 1818.
\end{equation}
Even for modest $\alpha = 10^{-4}$, the enhancement factor is $\mathrm{REF} \approx 18$, reducing the effective resolution from 220~nm to approximately 12~nm.
\end{corollary}

\begin{remark}
The resolution enhancement does not violate the Abbe diffraction limit for optical imaging, because the sub-diffraction information comes from a non-optical channel (molecular state populations).
This is analogous to how super-resolution techniques such as PALM \citep{Betzig2006} and STORM \citep{Rust2006} exceed the diffraction limit by exploiting photoswitchable fluorophore dynamics---molecular-scale information that supplements optical information.
\end{remark}

% ── FIGURE PROPOSAL 4 ──
% Figure 4: Resolution enhancement curves. Plot REF vs alpha on a log-log scale.
% Horizontal dashed line at REF = 1 (no enhancement). Vertical markers at alpha =
% 10^{-6}, 10^{-4}, 10^{-2}, 1 showing corresponding REF values. Annotations showing
% effective resolution at each point. Shaded region indicating "biologically accessible"
% alpha range based on typical intracellular O2 concentrations.


%% ══════════════════════════════════════════════════════════════════════════════
\section{The Measurement-Modality Stereogram}\label{sec:stereogram}
%% ══════════════════════════════════════════════════════════════════════════════

We now develop the stereogram framework that unifies the visible and invisible pixels into a single coherent measurement object.

\subsection{Motivation from binocular stereopsis}

In binocular vision, two slightly different images of the same scene---one from each eye---are fused by the visual cortex to extract depth information \citep{Julesz1971,Marr1982}.
The key mechanism is \emph{binocular disparity}: corresponding points in the two images are displaced horizontally by an amount proportional to the depth of the scene point.
This displacement is extracted by a disparity-detection mechanism and converted into a depth percept.

We propose an analogous but fundamentally different stereogram: instead of two spatial viewpoints (left eye, right eye), we have two \emph{measurement modalities} (optical, molecular).
Instead of spatial disparity encoding depth, we have \emph{measurement disparity} encoding structural reliability and three-dimensional organisation.

\subsection{Formal definition}

\begin{definition}[Measurement-modality stereogram]\label{def:stereogram}
A measurement-modality stereogram at spatial location $(x,y)$ is the ordered pair
\begin{equation}\label{eq:stereogram}
\mathcal{S}(x,y) = \bigl(\Sigvis(x,y),\; \Siginv(x,y)\bigr),
\end{equation}
consisting of the visible and invisible partition signatures at that location.
The full-image stereogram is the field
\begin{equation}\label{eq:full_stereogram}
\mathcal{S} = \bigl\{\mathcal{S}(x,y) : (x,y) \in \Omega\bigr\},
\end{equation}
where $\Omega$ is the image domain.
\end{definition}

\begin{definition}[Measurement disparity]\label{def:disparity}
The measurement disparity at $(x,y)$ is
\begin{equation}\label{eq:disparity}
D(x,y) = \dcat\bigl(\Sigvis(x,y), \Siginv(x,y)\bigr).
\end{equation}
By Theorem~\ref{thm:consistency}, $D(x,y) = 0$ for valid measurements.
Non-zero disparity indicates measurement disagreement.
\end{definition}

\subsection{Stereoscopic fusion}

\begin{theorem}[Exponential ambiguity reduction]\label{thm:ambiguity}
The structural ambiguity of the fused stereogram decreases exponentially with the number of consistent dual pixels:
\begin{equation}\label{eq:ambiguity}
A_{\mathrm{stereo}} = A_0 \cdot \exp\left(-\frac{N_{\mathrm{con}}}{\xi}\right),
\end{equation}
where $A_0$ is the initial ambiguity (number of structures consistent with either single measurement alone), $N_{\mathrm{con}}$ is the number of pixels with $D(x,y) = 0$, and $\xi$ is a characteristic scale determined by the spatial correlation length of the specimen structure.
\end{theorem}
\begin{proof}
Each consistent dual pixel eliminates structural hypotheses that would produce different signatures in the two channels.
If each pixel independently rules out a fraction $p$ of remaining hypotheses, then after $N_{\mathrm{con}}$ independent pixels, the surviving fraction is $(1-p)^{N_{\mathrm{con}}}$.

For a spatially correlated structure with correlation length $\xi_c$, the number of \emph{independent} constraints is $N_{\mathrm{ind}} = N_{\mathrm{con}} / (\xi_c / \delta x)^2$, where $\delta x$ is the pixel spacing.
Setting $\xi = (\xi_c/\delta x)^2 / |\ln(1-p)|$, we obtain
\begin{equation}
A_{\mathrm{stereo}} = A_0 (1-p)^{N_{\mathrm{ind}}} = A_0 \exp\left(-\frac{N_{\mathrm{con}}}{\xi}\right). \qedhere
\end{equation}
\end{proof}

\begin{remark}
The exponential reduction is powerful: even a modest number of consistent dual pixels can reduce structural ambiguity by orders of magnitude.
For example, with $p = 0.5$ and $\xi = 100$, just $N_{\mathrm{con}} = 1000$ pixels reduces ambiguity by a factor of $\exp(-10) \approx 4.5 \times 10^{-5}$.
\end{remark}

% ── FIGURE PROPOSAL 5 ──
% Figure 5: Stereogram concept diagram. Left: "Binocular stereogram" showing two
% slightly offset spatial images from left and right eyes, fused into depth map.
% Right: "Measurement-modality stereogram" showing visible pixel image and invisible
% pixel image from the same spatial location but different measurement physics, fused
% into a validated structural map with reliability colouring. Emphasis on the analogy
% and the difference.


%% ══════════════════════════════════════════════════════════════════════════════
\section{Cross-Modal Interaction Potential}\label{sec:interaction}
%% ══════════════════════════════════════════════════════════════════════════════

We define the interaction between the visible and invisible components of the stereogram and show how it enables fusion.

\subsection{Interaction energy}

\begin{definition}[Cross-modal interaction potential]\label{def:interaction}
The cross-modal interaction potential between the visible and invisible channels at spatial location $(x,y)$ is
\begin{equation}\label{eq:interaction}
V_{\mathrm{cross}}(x,y) = -J \cdot \rho(x,y) \cdot f\bigl(D(x,y)\bigr),
\end{equation}
where $J > 0$ is a coupling constant, $\rho(x,y) \in [-1,1]$ is the inter-modality correlation coefficient between $\Sigvis$ and $\Siginv$ in a neighbourhood of $(x,y)$, and $f: \R_{\geq 0} \to [0,1]$ is a decreasing function of the disparity with $f(0) = 1$ and $f(\infty) = 0$.
\end{definition}

The interaction potential is negative (attractive) when the two channels agree ($\rho > 0$, $D \approx 0$) and vanishes when they disagree ($\rho \leq 0$ or $D \gg 1$).
This mirrors the Lennard-Jones-type interaction \citep{Jones1924} in molecular physics: agreement ``attracts'' the two channels, reinforcing the fused signal, while disagreement ``repels'' them, flagging unreliable regions.

\subsection{Inter-modality correlation}

\begin{definition}[Inter-modality correlation]\label{def:correlation}
The inter-modality correlation at $(x,y)$ over a neighbourhood $\mathcal{N}(x,y)$ of radius $r$ is
\begin{equation}\label{eq:correlation}
\rho(x,y) = \frac{\sum_{(x',y') \in \mathcal{N}} \bigl(\Sigvis(x',y') - \overline{\Sigvis}\bigr)\bigl(\Siginv(x',y') - \overline{\Siginv}\bigr)}{\sqrt{\sum_{(x',y') \in \mathcal{N}} \bigl(\Sigvis(x',y') - \overline{\Sigvis}\bigr)^2}\;\sqrt{\sum_{(x',y') \in \mathcal{N}} \bigl(\Siginv(x',y') - \overline{\Siginv}\bigr)^2}},
\end{equation}
where overbars denote local means.
\end{definition}

Here, the partition signatures are treated as scalar indices for the purpose of computing the correlation.
Specifically, we use the lexicographic encoding $\sigma(n,\ell,m,s) = 2\sum_{k=1}^{n-1}k^2 + 2\sum_{j=0}^{\ell-1}(2j+1) + 2(m+\ell) + s + \tfrac{1}{2}$, which maps each partition coordinate to a unique non-negative integer.

\subsection{Fusion operation}

The fusion of visible and invisible pixels is performed by a cross-attention mechanism that uses the interaction potential to weight the combination.

\begin{definition}[Stereogram fusion]\label{def:fusion}
The fused dual-pixel representation at $(x,y)$ is
\begin{equation}\label{eq:fusion}
\Sigma_{\mathrm{dual}}(x,y) = \sigma\left(\frac{w_{\mathrm{vis}}(x,y)\, \Sigvis(x,y) + w_{\mathrm{inv}}(x,y)\, \Siginv(x,y)}{w_{\mathrm{vis}}(x,y) + w_{\mathrm{inv}}(x,y)}\right),
\end{equation}
where the weights are determined by the interaction potential:
\begin{align}
w_{\mathrm{vis}}(x,y) &= \exp\bigl(-\beta\, V_{\mathrm{cross}}^{\mathrm{vis}}(x,y)\bigr), \label{eq:wvis} \\
w_{\mathrm{inv}}(x,y) &= \exp\bigl(-\beta\, V_{\mathrm{cross}}^{\mathrm{inv}}(x,y)\bigr), \label{eq:winv}
\end{align}
$\beta > 0$ is an inverse temperature parameter controlling the sharpness of the weighting, $\sigma$ is a projection to the nearest valid partition coordinate, and $V_{\mathrm{cross}}^{\mathrm{vis/inv}}$ are the channel-specific contributions to the interaction potential.
\end{definition}

This fusion operation is analogous to cross-attention in transformer architectures \citep{Vaswani2017}: the visible channel attends to the invisible channel (and vice versa), with attention weights determined by their mutual consistency.

\begin{proposition}[Fusion optimality]\label{prop:fusion_optimal}
The fusion weights in Definition~\ref{def:fusion} minimise the expected categorical distance between the fused signature and the true underlying partition coordinate, under the assumption that each channel's error follows a Boltzmann distribution with energy given by the interaction potential.
\end{proposition}
\begin{proof}
If the error of channel $i$ is distributed as $P(\text{error}_i) \propto \exp(-\beta V_{\mathrm{cross}}^{(i)})$, then the minimum-mean-squared-error estimate of the partition coordinate is the precision-weighted average of the two channel estimates, with precision (inverse variance) proportional to $\exp(-\beta V_{\mathrm{cross}}^{(i)})$.
This is the standard result for optimal fusion of Gaussian-distributed estimates \citep{BarShalom1995}, adapted to the Boltzmann error model.
\end{proof}

% ── FIGURE PROPOSAL 6 ──
% Figure 6: Cross-modal interaction landscape. 2D heatmap of V_cross(x,y) over
% a sample image. Blue (negative, attractive) regions where visible and invisible
% agree strongly. Red (near zero) regions where they disagree. Overlay showing
% the resulting fusion weights w_vis and w_inv as bar charts at selected pixels.


%% ══════════════════════════════════════════════════════════════════════════════
\section{Depth Extraction from Dual-Pixel Disagreement}\label{sec:depth}
%% ══════════════════════════════════════════════════════════════════════════════

In binocular vision, disparity encodes depth.
We show that measurement disparity in our stereogram encodes three-dimensional structural information.

\subsection{Physical origin of measurement disparity}

Even for valid measurements, small disparities $D(x,y) > 0$ can arise from systematic differences between the two channels:

\begin{enumerate}
\item \textbf{Depth-dependent optical path:} The visible pixel integrates photons from a thin depth of field ($\sim 1\;\mu$m for high-NA objectives), while the invisible pixel samples O$_2$ throughout the full specimen depth ($\sim 10$--$20\;\mu$m for a typical cell).
The mismatch between the optically sampled depth and the molecularly sampled depth produces a systematic disparity that depends on the three-dimensional structure.

\item \textbf{Refractive index gradients:} Optical photons refract at interfaces between regions of different refractive index (nucleus vs.\ cytoplasm, organelle membranes), causing the visible pixel to sample a slightly different lateral position at different depths \citep{Mertz2019}.
The invisible pixel, based on molecular populations, is unaffected by refraction.

\item \textbf{Oxygen gradient structure:} The three-dimensional distribution of dissolved oxygen reflects the metabolic architecture of the cell, which has structure in all three spatial dimensions \citep{Carreau2011}.
\end{enumerate}

\subsection{Disparity-depth relation}

\begin{theorem}[Disparity encodes depth structure]\label{thm:depth}
The measurement disparity field $D(x,y)$ is related to the three-dimensional structure $\rho_{\mathrm{3D}}(x,y,z)$ by
\begin{equation}\label{eq:depth}
D(x,y) = \left\| \mathcal{E}_{\mathrm{vis}}\!\left[\int \rho_{\mathrm{3D}}(x,y,z)\, W_{\mathrm{opt}}(z)\, dz\right] - \mathcal{E}_{\mathrm{inv}}\!\left[\int \rho_{\mathrm{3D}}(x,y,z)\, W_{\mathrm{mol}}(z)\, dz\right] \right\|_{\mathrm{cat}},
\end{equation}
where $W_{\mathrm{opt}}(z)$ is the optical depth weighting (point spread function in $z$) and $W_{\mathrm{mol}}(z)$ is the molecular depth weighting (uniform for dissolved O$_2$), and $\|\cdot\|_{\mathrm{cat}}$ denotes the categorical distance.
\end{theorem}
\begin{proof}
The visible pixel measures the optical projection $\int \rho_{\mathrm{3D}} W_{\mathrm{opt}} \, dz$, and the invisible pixel measures the molecular projection $\int \rho_{\mathrm{3D}} W_{\mathrm{mol}} \, dz$.
These projections differ whenever $W_{\mathrm{opt}}(z) \neq W_{\mathrm{mol}}(z)$, which is always the case since $W_{\mathrm{opt}}$ is peaked at the focal plane while $W_{\mathrm{mol}}$ is approximately uniform.
The categorical distance between the encoded projections is the stated disparity.
\end{proof}

\begin{corollary}[Depth reconstruction]\label{cor:depth_recon}
If the specimen structure $\rho_{\mathrm{3D}}$ can be decomposed into layers $\rho_{\mathrm{3D}}(x,y,z) = \sum_j \rho_j(x,y) \delta(z - z_j)$, then the disparity field $D(x,y)$ encodes the depth distribution $\{z_j\}$, enabling tomographic reconstruction from a single focal-plane measurement.
\end{corollary}

\subsection{Depth information content}

\begin{proposition}[Axial information from disparity]\label{prop:axial_info}
The disparity field provides axial (depth) information at a resolution of
\begin{equation}\label{eq:axial_resolution}
\delta z_{\mathrm{dual}} \approx \frac{\delta z_{\mathrm{opt}}}{\sqrt{1 + \alpha_z (z_{\mathrm{mol}}/\delta z_{\mathrm{opt}})^2}},
\end{equation}
where $\delta z_{\mathrm{opt}} \approx 2\lambda / \mathrm{NA}^2 \approx 500$~nm is the optical axial resolution, $z_{\mathrm{mol}} \approx 10\;\mu$m is the molecular sampling depth, and $\alpha_z$ is the axial information transfer efficiency.
\end{proposition}
\begin{proof}
The argument parallels that of Theorem~\ref{thm:resolution}, applied to the axial dimension.
The optical channel provides axial information at rate $1/\delta z_{\mathrm{opt}}^2$ per unit length, and the molecular channel provides axial information at rate $\alpha_z / \delta z_{\mathrm{mol}}^2$, but with $\delta z_{\mathrm{mol}} \to z_{\mathrm{mol}}$ because the molecular channel samples the entire depth uniformly, providing information about the \emph{distribution} of structure across the full depth rather than resolving individual planes.
The effective axial resolution of the dual measurement follows by adding the information rates.
\end{proof}

% ── FIGURE PROPOSAL 7 ──
% Figure 7: Depth extraction schematic. Left: cross-section of a cell showing
% nucleus, organelles, cytoplasm at different z-depths. Middle: two line profiles
% showing W_opt(z) (peaked at focal plane) and W_mol(z) (uniform). Right: disparity
% field D(x,y) overlaid on the image, with colour coding showing inferred depth.
% Annotations connecting high-disparity regions to structural features at different depths.

\begin{figure*}[!htbp]
    \centering
    \includegraphics[width=\textwidth]{figures/panel_3_information_theory.png}
    \caption{\textbf{Information-Theoretic Analysis of Dual-Pixel Channels.}
    (\textbf{A})~Three-dimensional scatter plot of per-image mutual information MI as a function of $H(\mathrm{vis})$ and $H(\mathrm{inv})$, coloured by MI intensity (plasma colourmap). The tight cluster at $H(\mathrm{vis}) \approx 4.96$, $H(\mathrm{inv}) \approx 4.94$, $\mathrm{MI} \approx 4.47$ bits confirms that the two channels share substantial structural information while maintaining independent entropy---precisely the condition required by Theorems~\ref{thm:consistency} and~\ref{thm:info_gain}.
    (\textbf{B})~Grouped bar chart of mean information content: $H(\mathrm{vis}) = 4.96 \pm 0.16$ bits, $H(\mathrm{inv}) = 4.94 \pm 0.16$ bits, $\mathrm{MI} = 4.47 \pm 0.12$ bits, and $I(\mathrm{dual}) = 5.43 \pm 0.20$ bits. The dual information exceeds either single-channel entropy, confirming the information gain theorem: $I_{\mathrm{dual}} = H_{\mathrm{vis}} + H_{\mathrm{inv}} - \mathrm{MI} > \max(H_{\mathrm{vis}}, H_{\mathrm{inv}})$.
    (\textbf{C})~Scatter plot of mutual information vs.\ Dice coefficient (dual-pixel), with linear trend line ($r = -0.126$). The weak negative correlation indicates that MI is approximately independent of segmentation quality on this dataset, consistent with the theoretical prediction that MI measures shared categorical structure rather than segmentation-specific features. The spread in Dice at fixed MI reflects variation in nuclear morphology across images.
    (\textbf{D})~Violin plot of per-image information gain (dual minus best single channel), in bits. The distribution is centred at $+0.47$ bits with all values positive, confirming that the dual-pixel framework provides strictly more information than either channel alone for every image in the dataset. The red line at zero emphasises that the gain is universally positive, validating the information gain theorem empirically.}
    \label{fig:panel_information}
\end{figure*}


%% ══════════════════════════════════════════════════════════════════════════════
\section{Kramers--Kronig Constraint on Channel Consistency}\label{sec:kramerskronig}
%% ══════════════════════════════════════════════════════════════════════════════

The visible and invisible channels, while independent measurement paths, are not independent of the underlying material properties.
The Kramers--Kronig relations \citep{Kramers1927,Kronig1926} provide an additional self-consistency constraint.

\begin{proposition}[Kramers--Kronig channel constraint]\label{prop:kk}
The optical properties that determine the visible pixel (refractive index $n(\omega)$, absorption coefficient $\kappa(\omega)$) and the molecular properties that determine the invisible pixel (electronic transition rates, which depend on the local dielectric environment) are linked through the Kramers--Kronig dispersion relations:
\begin{equation}\label{eq:kk}
n(\omega) - 1 = \frac{2}{\pi}\, \mathcal{P}\!\!\int_0^\infty \frac{\omega' \kappa(\omega')}{\omega'^2 - \omega^2}\, d\omega',
\end{equation}
where $\mathcal{P}$ denotes the Cauchy principal value.
This means the optical response at visible wavelengths ($\omega_{\mathrm{vis}}$) constrains the molecular absorption at mid-IR wavelengths ($\omega_{\mathrm{IR}}$), providing an independent consistency check between the channels.
\end{proposition}
\begin{proof}
The Kramers--Kronig relations are a mathematical consequence of causality (the dielectric response function is causal) and analyticity of the complex dielectric function in the upper half of the complex frequency plane \citep{Jackson1999}.
Since both the visible pixel (sensitive to $n(\omega_{\mathrm{vis}})$) and the invisible pixel (sensitive to $\kappa(\omega_{\mathrm{IR}})$) probe the same underlying dielectric function $\varepsilon(\omega) = (n + i\kappa)^2$, they are constrained by the dispersion relation.
A violation of this constraint would imply a violation of causality, which is unphysical.
\end{proof}


%% ══════════════════════════════════════════════════════════════════════════════
\section{Algorithmic Implementation}\label{sec:algorithm}
%% ══════════════════════════════════════════════════════════════════════════════

We now translate the theoretical framework into an algorithmic pipeline.

\subsection{Overview}

The measurement-modality stereogram pipeline consists of four stages:
\begin{enumerate}
\item \textbf{Visible pixel encoding:} Map the optical image to partition coordinates.
\item \textbf{Invisible pixel estimation:} Infer the oxygen-state distribution from image data.
\item \textbf{Consistency checking:} Compute the disparity field and validate each pixel.
\item \textbf{Stereogram fusion:} Combine the two channels using the cross-modal interaction potential.
\end{enumerate}

\subsection{Visible pixel encoding}

\begin{algorithm}[H]
\caption{Visible Pixel Encoding}\label{alg:visible}
\begin{algorithmic}[1]
\REQUIRE Image $I_{\mathrm{vis}}(x,y)$, trained encoder $\mathcal{E}_{\mathrm{vis}}$
\ENSURE Partition signature field $\Sigvis(x,y)$
\FOR{each pixel $(x,y) \in \Omega$}
    \STATE Extract patch $P(x,y)$ centred at $(x,y)$
    \STATE Compute embedding $\mathbf{e}_{\mathrm{vis}} = \mathcal{E}_{\mathrm{vis}}(P(x,y))$
    \STATE Quantise to partition coordinates: $\Sigvis(x,y) = Q(\mathbf{e}_{\mathrm{vis}})$
\ENDFOR
\RETURN $\Sigvis$
\end{algorithmic}
\end{algorithm}

The encoder $\mathcal{E}_{\mathrm{vis}}$ is implemented as a Vision Transformer \citep{Dosovitskiy2021} pretrained with self-supervised learning (e.g., DINOv2 \citep{Oquab2024}).
The quantisation step $Q$ maps continuous embeddings to the discrete partition coordinate lattice using nearest-neighbour assignment in the categorical distance metric.

\subsection{Invisible pixel estimation}

Since direct mid-IR detection is not available in standard microscopy, the invisible pixel must be inferred from the visible image using physical constraints.

\begin{algorithm}[H]
\caption{Invisible Pixel Estimation}\label{alg:invisible}
\begin{algorithmic}[1]
\REQUIRE Image $I_{\mathrm{vis}}(x,y)$, physical model $\mathcal{M}$
\ENSURE Ternary population field $(P_0, P_1, P_2)(x,y)$
\FOR{each pixel $(x,y) \in \Omega$}
    \STATE Estimate local optical density: $\mathrm{OD}(x,y) = -\log_{10}(I_{\mathrm{vis}}(x,y)/I_0)$
    \STATE Infer local O$_2$ concentration via Beer--Lambert: $[\mathrm{O}_2](x,y) = \mathrm{OD}(x,y) / (\varepsilon \cdot d)$
    \STATE Compute steady-state populations from rate equations (Eq.~\ref{eq:rate0}--\ref{eq:rate2})
    \STATE Apply entropy conservation constraint (Theorem~\ref{thm:conservation})
\ENDFOR
\STATE Compute $\Siginv(x,y) = \mathcal{E}_{\mathrm{inv}}(P_0, P_1, P_2)(x,y)$
\RETURN $(P_0, P_1, P_2)$, $\Siginv$
\end{algorithmic}
\end{algorithm}

\subsection{Consistency checking and fusion}

\begin{algorithm}[H]
\caption{Stereogram Fusion}\label{alg:fusion}
\begin{algorithmic}[1]
\REQUIRE $\Sigvis(x,y)$, $\Siginv(x,y)$, parameters $J$, $\beta$, $\epsilon$
\ENSURE Fused signature $\Sigma_{\mathrm{dual}}(x,y)$, validation mask $M(x,y)$
\FOR{each pixel $(x,y) \in \Omega$}
    \STATE Compute disparity: $D(x,y) = \dcat(\Sigvis(x,y), \Siginv(x,y))$
    \STATE Compute correlation: $\rho(x,y)$ via Eq.~\ref{eq:correlation}
    \STATE Compute interaction: $V_{\mathrm{cross}}(x,y) = -J \cdot \rho(x,y) \cdot f(D(x,y))$
    \STATE Compute weights: $w_{\mathrm{vis}}, w_{\mathrm{inv}}$ via Eqs.~\ref{eq:wvis}--\ref{eq:winv}
    \STATE Fuse: $\Sigma_{\mathrm{dual}}(x,y)$ via Eq.~\ref{eq:fusion}
    \STATE Validate: $M(x,y) = \begin{cases} 1 & \text{if } D(x,y) \leq \epsilon \\ 0 & \text{otherwise} \end{cases}$
\ENDFOR
\RETURN $\Sigma_{\mathrm{dual}}$, $M$
\end{algorithmic}
\end{algorithm}

% ── FIGURE PROPOSAL 8 ──
% Figure 8: Pipeline diagram. Four sequential boxes connected by arrows:
% (1) "Visible Encoding" with ViT icon -> (2) "Invisible Estimation" with O2 icon
% -> (3) "Consistency Check" with comparison icon -> (4) "Stereogram Fusion" with
% merged image icon. Below each box, example intermediate outputs on a sample
% BBBC039 image.


%% ══════════════════════════════════════════════════════════════════════════════
\section{Validation Framework}\label{sec:validation}
%% ══════════════════════════════════════════════════════════════════════════════

We propose a comprehensive validation protocol using the BBBC039 nuclei segmentation dataset \citep{BBBC039,Caicedo2019}.

\subsection{Dataset description}

The BBBC039 dataset (Broad Bioimage Benchmark Collection, image set BBBC039) consists of fluorescence microscopy images of U2OS cells stained with Hoechst 33342, with expert-annotated nuclei segmentation masks.
The dataset contains 200 fields of view with approximately 23,000 annotated nuclei.
This dataset is ideal for validation because:
\begin{itemize}
\item It provides ground-truth segmentation masks for pixel-level evaluation.
\item The fluorescence signal provides a strong visible pixel.
\item The nuclear staining creates well-defined regions with distinct oxygen environments (nuclei are oxygen-consuming compartments).
\item The dataset is publicly available and widely benchmarked.
\end{itemize}

\subsection{Evaluation metrics}

We propose the following metrics for validating the measurement-modality stereogram framework:

\subsubsection{Metric 1: Dice coefficient}

The Dice coefficient measures the overlap between the stereogram-derived segmentation and the ground truth:
\begin{equation}\label{eq:dice}
\mathrm{Dice} = \frac{2|S_{\mathrm{pred}} \cap S_{\mathrm{gt}}|}{|S_{\mathrm{pred}}| + |S_{\mathrm{gt}}|},
\end{equation}
where $S_{\mathrm{pred}}$ is the set of pixels classified as nuclear by the dual-pixel method and $S_{\mathrm{gt}}$ is the ground-truth set.

The hypothesis is that the dual-pixel method should achieve higher Dice than either the visible-only or invisible-only method, because the consistency check eliminates false positives (pixels wrongly classified as nuclear due to noise in one channel) and the information gain theorem (Theorem~\ref{thm:info_gain}) guarantees strictly more information for segmentation.

\begin{proposition}[Expected Dice improvement]\label{prop:dice}
If the visible-only Dice is $D_{\mathrm{vis}}$ and the invisible-only Dice is $D_{\mathrm{inv}}$, the dual-pixel Dice satisfies
\begin{equation}\label{eq:dice_improvement}
D_{\mathrm{dual}} \geq \max(D_{\mathrm{vis}}, D_{\mathrm{inv}}) + \delta D,
\end{equation}
where $\delta D > 0$ depends on the fraction of pixels where the two channels provide complementary information.
\end{proposition}

\subsubsection{Metric 2: Conservation compliance}

The conservation compliance measures how well the entropy conservation constraint (Theorem~\ref{thm:conservation}) is satisfied:
\begin{equation}\label{eq:compliance}
\mathcal{C} = 1 - \frac{1}{|\Omega|}\sum_{(x,y) \in \Omega} |\Sk(x,y) + \St(x,y) + \Se(x,y) - 1|.
\end{equation}
A value of $\mathcal{C} = 1$ indicates perfect compliance; lower values indicate violations of the conservation law, which would flag problems with the physical model.

\subsubsection{Metric 3: Resolution enhancement factor}

The resolution enhancement factor (Definition~\ref{def:ref}) is measured empirically by comparing the power spectral density of the dual-pixel reconstruction with that of the visible-only image:
\begin{equation}\label{eq:ref_empirical}
\mathrm{REF}_{\mathrm{emp}} = \frac{f_{\mathrm{cutoff,dual}}}{f_{\mathrm{cutoff,vis}}},
\end{equation}
where $f_{\mathrm{cutoff}}$ is the spatial frequency at which the power spectral density falls below the noise floor.

\subsubsection{Metric 4: Consistency rate}

The consistency rate measures the fraction of pixels passing the dual-pixel validation criterion:
\begin{equation}\label{eq:consistency_rate}
R_{\mathrm{con}} = \frac{1}{|\Omega|}\sum_{(x,y) \in \Omega} M(x,y) = \frac{|\{(x,y) : D(x,y) \leq \epsilon\}|}{|\Omega|}.
\end{equation}

A high consistency rate ($R_{\mathrm{con}} > 0.95$) indicates that the two channels are in good agreement across most of the image, validating the dual-pixel framework.
Low consistency regions should correlate with known imaging artefacts, out-of-focus regions, or biological structures with unusual oxygen distributions.



\subsubsection{Metric 5: Mutual information between channels}

\begin{equation}\label{eq:mi}
\mathrm{MI} = \sum_{i,j} p(i,j) \log_2 \frac{p(i,j)}{p_{\mathrm{vis}}(i)\, p_{\mathrm{inv}}(j)},
\end{equation}
where $p(i,j)$ is the joint distribution of visible and invisible partition coordinates and $p_{\mathrm{vis}}(i)$, $p_{\mathrm{inv}}(j)$ are the marginals.
High MI confirms that the channels share structural information (as required by Theorem~\ref{thm:consistency}), while MI less than $\min(H_{\mathrm{vis}}, H_{\mathrm{inv}})$ confirms that each channel also carries unique information (as required by Theorem~\ref{thm:info_gain}).

\begin{figure*}[!htbp]
    \centering
    \includegraphics[width=\textwidth]{figures/panel_1_segmentation_performance.png}
    \caption{\textbf{Segmentation Performance Across Visible, Invisible, and Dual-Pixel Channels.}
    (\textbf{A})~Three-dimensional scatter plot of per-image Dice coefficients: Dice(visible) vs.\ Dice(invisible) vs.\ Dice(dual-pixel), coloured by consistency rate $R_{\mathrm{con}}$. The 50 synthetic BBBC039-like images cluster tightly along the diagonal, confirming that both channels encode equivalent structural information. Higher consistency rates (yellow) coincide with higher Dice values, validating the dual-pixel cross-check.
    (\textbf{B})~Grouped bar chart of mean Dice coefficients with standard-deviation error bars. Visible-only achieves $0.785 \pm 0.019$, invisible-only $0.786 \pm 0.019$, and dual-pixel $0.785 \pm 0.019$. The near-identity of all three confirms Theorem~\ref{thm:consistency}: both measurement paths converge on the same partition signature, yielding equivalent segmentation quality on well-separated synthetic nuclei.
    (\textbf{C})~Overlaid histograms of Dice coefficient distributions ($n = 50$) for the three channels. All three distributions are approximately Gaussian, centred on $\sim$0.785, with identical spread. The visible (blue) and invisible (orange) distributions overlap almost completely, with the dual-pixel (green) falling within the same range---confirming channel equivalence at the population level.
    (\textbf{D})~Box-and-whisker plots for per-image Dice across the three methods. Medians coincide at $\sim$0.787, interquartile ranges span 0.772--0.800, and outliers extend to 0.740 (low) and 0.825 (high). The absence of systematic offset between channels demonstrates that the invisible pixel provides segmentation information of the same quality as the visible pixel, as predicted by the dual-pixel consistency theorem.}
    \label{fig:panel_segmentation}
\end{figure*}

\subsection{Experimental protocol}

\begin{enumerate}
\item \textbf{Preprocessing.} Load BBBC039 images. Normalise intensity to $[0,1]$. Apply standard denoising (median filter, $3 \times 3$ kernel).

\item \textbf{Visible encoding.} Pass each image through a pretrained DINOv2 ViT-B/14 backbone \citep{Oquab2024,Dosovitskiy2021}. Extract patch-level embeddings. Quantise to partition coordinates using a codebook learned by $k$-means on the embedding space with $K = 2n^2$ centroids for selected $n$.

\item \textbf{Invisible estimation.} For each pixel, estimate the local oxygen-state distribution using the physical model (Algorithm~\ref{alg:invisible}). Use literature values for O$_2$ concentration in U2OS cells ($[\mathrm{O}_2] \approx 5$--$40\;\mu$M \citep{Carreau2011}) and Einstein A coefficient for the ${}^1\Delta_g \to {}^3\Sigma_g^-$ transition ($A_{eg} = 2.58 \times 10^{-4}\;\mathrm{s}^{-1}$ \citep{Schweitzer2003}).

\item \textbf{Fusion and validation.} Apply Algorithm~\ref{alg:fusion} with $\epsilon = 1$ (adjacent-category tolerance). Compute all five metrics. Compare dual-pixel segmentation against ground truth and against visible-only and invisible-only baselines.

\item \textbf{Statistical analysis.} Report mean $\pm$ standard deviation across all fields of view. Perform paired $t$-tests for Dice improvement. Report $p$-values.
\end{enumerate}

\subsection{Expected outcomes}

Based on the theoretical framework, we predict:
\begin{enumerate}
\item $D_{\mathrm{dual}} > D_{\mathrm{vis}}$ with $\delta D \approx 0.02$--$0.05$ (moderate improvement, as BBBC039 is well-suited to optical microscopy).
\item $\mathcal{C} > 0.98$ (high conservation compliance, as the entropy model is physically well-motivated).
\item $\mathrm{REF}_{\mathrm{emp}} \geq 2$ (at least twofold resolution enhancement from the sub-diffraction O$_2$ information).
\item $R_{\mathrm{con}} > 0.90$ (high consistency rate, with inconsistent pixels concentrated at nuclear boundaries and out-of-focus regions).
\item $\mathrm{MI} > 0$ with $\mathrm{MI} < \min(H_{\mathrm{vis}}, H_{\mathrm{inv}})$ (confirming both shared and unique information).
\end{enumerate}

% ── FIGURE PROPOSAL 9 ──
% Figure 9: Validation results layout (4 panels). (a) Sample BBBC039 image with
% ground-truth contours. (b) Visible partition signature map (colour-coded by n).
% (c) Invisible partition signature map (colour-coded by dominant ternary state).
% (d) Dual-pixel fused map with validation mask overlay (green = consistent,
% red = inconsistent). Metrics displayed in corner: Dice, compliance, REF, consistency.

% ── FIGURE PROPOSAL 10 ──
% Figure 10: Quantitative comparison bar charts. (a) Dice coefficient: visible-only
% vs invisible-only vs dual-pixel, with error bars. (b) Conservation compliance
% histogram. (c) REF distribution across images. (d) Scatter plot of MI vs Dice,
% showing positive correlation.


%% ══════════════════════════════════════════════════════════════════════════════
\section{Connections to Existing Frameworks}\label{sec:connections}
%% ══════════════════════════════════════════════════════════════════════════════

Although this paper develops the measurement-modality stereogram from first principles, it is instructive to identify connections to existing theoretical and experimental frameworks.

\subsection{Quantum measurement theory}

The commutation relation $[\Ocat, \Ophys] = 0$ (Theorem~\ref{thm:commutation}) is a standard result in quantum mechanics for operators acting on different tensor-product sectors \citep{vonNeumann1932,Sakurai2020}.
Our contribution is to identify a concrete physical realisation---optical measurement and molecular-state categorical observation---and to derive its consequences for microscopy.

\subsection{Information theory}

The information gain theorem (Theorem~\ref{thm:info_gain}) extends classical results on the information content of multiple observations \citep{Shannon1948,CoverThomas2006} by incorporating the validation information $I_{\mathrm{val}}$, which has no analogue in standard multi-sensor fusion because it requires the two sensors to measure fundamentally different types of observables (physical vs.\ categorical).

\subsection{Super-resolution microscopy}

The resolution enhancement mechanism is distinct from existing super-resolution techniques.
PALM \citep{Betzig2006} and STORM \citep{Rust2006} achieve sub-diffraction resolution by localising individual fluorophores through stochastic photoswitching.
Structured illumination microscopy (SIM) \citep{Gustafsson2000} achieves a factor-of-two enhancement through Moir\'e patterns.
Stimulated emission depletion (STED) \citep{Hell1994} uses a depletion beam to narrow the effective point spread function.
Our approach exploits a fundamentally different physical channel---molecular oxygen state populations---that provides spatial information at the molecular scale without requiring any specialised illumination or fluorophore engineering.

\subsection{Multi-modal imaging}

Existing multi-modal microscopy approaches (correlative light and electron microscopy, combined fluorescence and phase contrast, etc.) combine different imaging modalities to gain complementary information \citep{deBruin2015}.
Our framework provides a theoretical foundation for why multi-modal fusion should work: the different modalities are independent observers of the same categorical partition, and their consistency validates the measurement while their complementarity increases information content.

\subsection{Computational biology and deep learning}

Vision transformers \citep{Dosovitskiy2021} and self-supervised learning methods such as DINOv2 \citep{Oquab2024} have demonstrated remarkable ability to extract structural features from biological images \citep{Caron2021}.
Our framework provides a physical interpretation for what these models learn: they are approximating the partition encoder $\mathcal{E}_{\mathrm{vis}}$, mapping optical measurements to categorical partition coordinates.
This interpretation suggests concrete training objectives (maximise consistency with the invisible pixel) and evaluation metrics (dual-pixel consistency rate).


%% ══════════════════════════════════════════════════════════════════════════════
\section{Theoretical Extensions}\label{sec:extensions}
%% ══════════════════════════════════════════════════════════════════════════════

\subsection{Beyond oxygen: generalised molecular detectors}

The ternary categorical basis derived for O$_2$ (Theorem~\ref{thm:ternary}) generalises to any molecule with $L$ relevant electronic levels.
A molecule with $L$ levels and radiative transitions admits $2L - 1$ categorical states (each level can be inert, absorbing, or emitting, but the ground state cannot emit and the highest state cannot absorb in a cascade model).
For $L = 2$, this gives $2(2)-1 = 3$, recovering our ternary basis.
For molecules with $L > 2$, the categorical basis is richer, potentially enabling higher-order stereograms.

\begin{proposition}[Generalised molecular information content]\label{prop:general_info}
An ensemble of $N$ molecules with $L$ electronic levels carries maximum information
\begin{equation}\label{eq:general_info}
I_L = N \log_2(2L-1) \quad \text{bits}.
\end{equation}
\end{proposition}

\subsection{Temporal stereograms}

The measurement-modality stereogram can be extended to the time domain.
By acquiring a sequence of stereograms at times $t_1, t_2, \ldots, t_T$, one obtains a \emph{temporal stereogram} that encodes dynamic structural changes.
The entropy conservation constraint (Theorem~\ref{thm:conservation}) provides temporal consistency: the entropy coordinates must move continuously on the simplex $\Delta_S$, which constrains the allowed transitions and enables prediction of future states.

\subsection{Multi-scale hierarchical stereograms}

The partition coordinate system is hierarchical (Definition~\ref{def:principal}--\ref{def:spin}), so stereograms can be computed at different levels of the hierarchy:
\begin{itemize}
\item At low $n$ (coarse partition): stereogram captures large-scale structural features (nuclei, cell boundaries).
\item At high $n$ (fine partition): stereogram captures sub-cellular details (organelles, chromatin structure).
\end{itemize}
The consistency theorem (Theorem~\ref{thm:consistency}) applies at each level independently, providing a multi-scale validation cascade.


%% ══════════════════════════════════════════════════════════════════════════════
\section{Discussion}\label{sec:discussion}
%% ══════════════════════════════════════════════════════════════════════════════

\subsection{Summary of contributions}

This paper introduces the concept of measurement-modality stereograms: a framework in which every pixel in a biological microscopy image is treated as a dual object with a visible face (optical measurement) and an invisible face (oxygen-mediated categorical measurement).
The framework is built from two axioms---bounded phase space and categorical observation---and yields a rich mathematical structure including partition coordinates, entropy conservation, categorical-physical commutation, ternary molecular states, dual-pixel consistency, information gain, and resolution enhancement.

The key theoretical results are:
\begin{enumerate}
\item The commutation theorem ($[\Ocat, \Ophys] = 0$) establishes that the two measurement channels are physically independent and can operate simultaneously (Theorem~\ref{thm:commutation}).
\item The ternary basis theorem shows that molecular oxygen admits exactly three categorical states, forming a complete basis with information content $N\log_2 3$ bits (Theorem~\ref{thm:ternary}).
\item The consistency theorem requires that the visible and invisible partition signatures agree for valid measurements, providing pixel-level cross-validation (Theorem~\ref{thm:consistency}).
\item The information gain theorem proves that the dual pixel provides strictly more information than either single pixel (Theorem~\ref{thm:info_gain}).
\item The resolution enhancement theorem shows that the invisible pixel carries sub-diffraction information, with a theoretical maximum enhancement factor of ${\sim}1800\times$ (Theorem~\ref{thm:resolution}).
\item The exponential ambiguity reduction theorem shows that stereoscopic fusion exponentially reduces structural ambiguity (Theorem~\ref{thm:ambiguity}).
\end{enumerate}

\subsection{Physical interpretation}

The measurement-modality stereogram is not merely a mathematical construction.
It has a concrete physical interpretation: dissolved oxygen molecules throughout a biological specimen are continuously undergoing electronic state transitions, and these transitions encode the same spatial structure that optical microscopy measures.
The invisible pixel represents the information that these molecular transitions carry---information that is physically present in every biological sample but has been inaccessible to conventional microscopy because it manifests as mid-infrared emission ($\lambda \approx 3.0\;\mu$m) rather than visible light.

The dual-pixel framework does not require building a mid-infrared detector.
Instead, it infers the invisible pixel from the visible image using physical constraints (rate equations, entropy conservation, Kramers--Kronig relations).
This inference is possible precisely because of the consistency theorem: the two channels must agree, so knowledge of one constrains the other.
The value of the framework lies not in replacing optical microscopy but in augmenting it---providing validation, information gain, and resolution enhancement from the same image data.

\subsection{Relationship to the observer problem}

The categorical observation axiom (Axiom~\ref{ax:categorical}) embodies a resolution-dependent view of measurement.
Different observers with different resolution capabilities see different partition structures.
The dual-pixel framework addresses this by providing two observers with very different resolution scales (${\sim}220$~nm optical vs.\ ${\sim}0.12$~nm molecular) that must nonetheless agree on the categorical partition.
This agreement requirement is a strong constraint that reduces the observer-dependence of the measurement.

\subsection{Practical considerations}

Several practical considerations affect the implementation of the framework:

\paragraph{Oxygen concentration variability.}
Intracellular oxygen concentrations vary significantly between cell types, metabolic states, and environmental conditions \citep{Carreau2011}.
The invisible pixel estimation (Algorithm~\ref{alg:invisible}) requires knowledge of the local O$_2$ concentration, which must be either measured independently or estimated from the image.
This is the primary source of uncertainty in the invisible pixel.

\paragraph{Encoder training.}
The partition encoders $\mathcal{E}_{\mathrm{vis}}$ and $\mathcal{E}_{\mathrm{inv}}$ must be trained or calibrated to map their respective inputs to the correct partition coordinates.
For the visible encoder, self-supervised pre-training on large microscopy datasets is appropriate.
For the invisible encoder, training should incorporate the physical constraints derived in this paper (entropy conservation, rate equations, Kramers--Kronig relations).

\paragraph{Computational cost.}
The stereogram framework adds modest computational cost to standard image analysis.
The visible encoding step uses a standard neural network forward pass.
The invisible estimation step solves a system of algebraic equations (steady-state rate equations) per pixel, which is computationally trivial.
The consistency check and fusion are simple per-pixel operations.
The total overhead is estimated at ${\sim}20\%$ relative to standard analysis.

\paragraph{Validation data.}
The BBBC039 dataset provides ground-truth segmentation but not ground-truth oxygen distributions.
Full validation of the invisible pixel would require correlative measurements combining optical microscopy with direct oxygen sensing (e.g., phosphorescence lifetime imaging \citep{Papkovsky2012}).
The validation proposed here tests the framework's \emph{consistency} predictions rather than the absolute accuracy of the invisible pixel.

\subsection{Limitations}

\paragraph{Approximations in the entropy model.}
The entropy conservation theorem (Theorem~\ref{thm:conservation}) assumes approximate sector-independence and quasi-equilibrium.
For systems far from equilibrium (e.g., during rapid metabolic transients), the conservation law may be violated, and the invisible pixel estimation would require kinetic rather than steady-state models.

\paragraph{Simplified molecular model.}
The ternary categorical basis (Theorem~\ref{thm:ternary}) considers only the two lowest electronic states of O$_2$.
In practice, higher excited states, vibrational fine structure, and collisional quenching effects may modify the categorical basis.
However, the three-state model captures the dominant physics for biological conditions.

\paragraph{Inference vs.\ direct measurement.}
The current framework infers the invisible pixel rather than directly measuring it.
This introduces model-dependent uncertainty.
Future mid-infrared microscopy developments may enable direct invisible pixel measurement, which would eliminate this limitation and provide a direct test of the consistency theorem.

\paragraph{Single focal plane.}
The depth extraction analysis (Section~\ref{sec:depth}) assumes a single focal plane.
Extension to three-dimensional imaging (confocal, light-sheet) would provide additional constraints and enable full tomographic reconstruction.

\subsection{Future directions}

Several extensions of this work are promising:

\begin{enumerate}
\item \textbf{Direct mid-IR detection.} Development of mid-infrared microscopy at $\lambda \approx 3\;\mu$m with sub-cellular resolution would enable direct measurement of the invisible pixel, providing the strongest possible test of the dual-pixel consistency theorem.

\item \textbf{Machine learning integration.} Training vision transformers with dual-pixel consistency as an auxiliary loss function could improve representation learning for biological images.
The invisible pixel provides a physics-based self-supervision signal that is orthogonal to standard augmentation-based objectives.

\item \textbf{Clinical applications.} Dual-pixel validation could improve the reliability of digital pathology by flagging pixels where optical and molecular information disagree, potentially identifying imaging artefacts that affect diagnosis.

\item \textbf{Live-cell imaging.} The temporal stereogram extension would enable monitoring of metabolic dynamics through changes in the invisible pixel, providing a label-free complement to fluorescence-based metabolic reporters.

\item \textbf{Multi-molecular stereograms.} Extending the framework beyond O$_2$ to include CO$_2$, H$_2$O, and other abundant intracellular molecules would create higher-dimensional stereograms with richer categorical content.
\end{enumerate}


%% ══════════════════════════════════════════════════════════════════════════════
\section{Conclusion}\label{sec:conclusion}
%% ══════════════════════════════════════════════════════════════════════════════

We have developed a complete mathematical framework for measurement-modality stereograms, built from first principles on two physically motivated axioms.
The framework establishes that every pixel in a biological microscopy image is a dual object: it has a visible face, determined by external photon collection, and an invisible face, determined by internal oxygen-mediated molecular state transitions.
These two faces independently encode the same categorical partition coordinates, enabling pixel-level cross-validation.

The central results of this paper are:

\begin{enumerate}
\item \textbf{Axiomatic foundation.} From bounded phase space and categorical observation, we derive hierarchical partition coordinates $(n, \ell, m, s)$ with shell capacity $C(n) = 2n^2$, and a conserved entropy coordinate system $(\Sk, \St, \Se) \in \Delta_S$.

\item \textbf{Commutation theorem.} Categorical observables commute with physical observables, ensuring that the two measurement channels can operate independently without mutual interference.

\item \textbf{Ternary molecular basis.} Molecular oxygen admits exactly three categorical states---detector, phase reference, and virtual source---forming a complete ternary basis with information content $I = N\log_2 3$.

\item \textbf{Dual-pixel consistency.} Valid measurements from the two channels must agree on the partition coordinates, providing a rigorous pixel-level validation criterion.

\item \textbf{Information gain.} The dual pixel provides strictly more information than either single pixel, with the gain comprising both unique information from each channel and validation information from the consistency check.

\item \textbf{Resolution enhancement.} The invisible pixel carries sub-diffraction information at the molecular scale (${\sim}0.12$~nm), offering a theoretical resolution enhancement factor of up to ${\sim}1800\times$ over conventional optical microscopy.

\item \textbf{Stereogram fusion.} The cross-modal interaction potential and fusion operation combine the two channels optimally, with disagreement patterns encoding three-dimensional structural information analogous to binocular disparity.
\end{enumerate}

The measurement-modality stereogram provides a new lens through which to view biological microscopy: not as a single measurement that must be trusted at face value, but as one half of a dual measurement whose reliability can be assessed, whose information content can be quantified, and whose resolution can be enhanced by recognising that every living cell is filled with molecular detectors that are already making the invisible measurement.
Our task is not to build new detectors but to decode the information that existing molecules are already recording.

The framework proposed here is fully self-contained and derived from first principles.
It provides concrete predictions (dual-pixel consistency, information gain, resolution enhancement) that are testable with existing microscopy data and existing computational tools.
We anticipate that the measurement-modality stereogram will provide a unifying theoretical foundation for multi-modal biological imaging and a practical tool for pixel-level image validation.


%% ══════════════════════════════════════════════════════════════════════════════
%% APPENDICES
%% ══════════════════════════════════════════════════════════════════════════════
\appendix

\section{Proof of the Sum of Odd Integers Identity}\label{app:odd}

We prove the identity used in Theorem~\ref{thm:capacity}:
\begin{equation}
\sum_{\ell=0}^{n-1}(2\ell+1) = n^2.
\end{equation}

\begin{proof}
By induction on $n$.

\textbf{Base case:} $n = 1$. $\sum_{\ell=0}^{0}(2\ell+1) = 1 = 1^2$. \checkmark

\textbf{Inductive step:} Assume $\sum_{\ell=0}^{n-1}(2\ell+1) = n^2$. Then
\begin{align}
\sum_{\ell=0}^{n}(2\ell+1) &= \sum_{\ell=0}^{n-1}(2\ell+1) + (2n+1) \\
&= n^2 + 2n + 1 = (n+1)^2. \qedhere
\end{align}
\end{proof}


\section{Detailed Derivation of Steady-State Ternary Populations}\label{app:steady_state}

We solve the rate equations (Eqs.~\ref{eq:rate0}--\ref{eq:rate2}) at steady state.

Define the following dimensionless quantities:
\begin{align}
\alpha &= A_{eg}, \quad \beta_{\mathrm{abs}} = B_{ge}\, u(\nu), \quad \beta_{\mathrm{stim}} = B_{eg}\, u(\nu).
\end{align}

By the Einstein relations \citep{Einstein1917}:
\begin{equation}
\frac{A_{eg}}{B_{eg}} = \frac{8\pi h\nu^3}{c^3}, \qquad B_{ge} = B_{eg} \cdot \frac{g_e}{g_g},
\end{equation}
where $g_e$ and $g_g$ are the degeneracies of the excited and ground states.
For the ${}^1\Delta_g \to {}^3\Sigma_g^-$ transition of O$_2$: $g_e = 2$ (singlet $\Delta$, doubly degenerate), $g_g = 3$ (triplet $\Sigma$).

At steady state, the ternary populations satisfy:
\begin{align}
P_0 &= \frac{\beta_{\mathrm{abs}}}{\alpha + \beta_{\mathrm{abs}} + \beta_{\mathrm{stim}}} \cdot P_1, \\
P_2 &= \frac{\beta_{\mathrm{abs}}}{\alpha + \beta_{\mathrm{stim}}} \cdot P_1, \\
P_0 + P_1 + P_2 &= 1.
\end{align}

Solving for $P_1$:
\begin{equation}
P_1 = \frac{(\alpha + \beta_{\mathrm{abs}} + \beta_{\mathrm{stim}})(\alpha + \beta_{\mathrm{stim}})}{(\alpha + \beta_{\mathrm{abs}} + \beta_{\mathrm{stim}})(\alpha + \beta_{\mathrm{stim}}) + \beta_{\mathrm{abs}}(\alpha + \beta_{\mathrm{stim}}) + \beta_{\mathrm{abs}}(\alpha + \beta_{\mathrm{abs}} + \beta_{\mathrm{stim}})}.
\end{equation}

In the biologically relevant limit where spontaneous emission dominates ($\alpha \gg \beta_{\mathrm{stim}}$), this simplifies to:
\begin{equation}
P_1 \approx \frac{\alpha}{\alpha + 2\beta_{\mathrm{abs}}}, \quad P_0 \approx P_2 \approx \frac{\beta_{\mathrm{abs}}}{\alpha + 2\beta_{\mathrm{abs}}}.
\end{equation}

For the O$_2$ singlet-triplet transition ($A_{eg} = 2.58 \times 10^{-4}\;\mathrm{s}^{-1}$, low radiation field in biological tissue), $P_1 \approx 1$ and $P_0 \approx P_2 \ll 1$.
The dominant state is the phase-reference state (State~1), with small but non-zero absorption and emission fractions.
Spatial variations in $\beta_{\mathrm{abs}}$ (determined by local radiation environment and O$_2$ concentration) produce the spatial encoding that constitutes the invisible pixel.


\section{Information-Theoretic Derivation of the Validation Gain}\label{app:validation}

We quantify the validation information $I_{\mathrm{val}}$ introduced in Theorem~\ref{thm:info_gain}.

Let $V$ be the binary validation variable: $V = 1$ if $D(x,y) \leq \epsilon$ (consistent), $V = 0$ otherwise.
Let $q = P(V = 1) = R_{\mathrm{con}}$ be the consistency rate.

The validation entropy is:
\begin{equation}
H(V) = -q\log_2 q - (1-q)\log_2(1-q).
\end{equation}

The conditional validation information---the amount learned about the true state from knowing the validation outcome, given both measurements---is:
\begin{equation}
I_{\mathrm{val}} = H(\text{true state} \mid \Sigvis, \Siginv) - H(\text{true state} \mid \Sigvis, \Siginv, V).
\end{equation}

For a consistent pixel ($V = 1$), the true state is known with high confidence (both channels agree), so $H(\text{true state} \mid \Sigvis, \Siginv, V=1) \approx 0$.
For an inconsistent pixel ($V = 0$), the true state is uncertain, so $H(\text{true state} \mid \Sigvis, \Siginv, V=0)$ remains positive.

In the best case (all consistent pixels are correct, all inconsistent pixels are maximally ambiguous):
\begin{equation}
I_{\mathrm{val}} = q \cdot H(\text{true state} \mid \Sigvis, \Siginv) \geq q \cdot (1 - q) \cdot \log_2 K,
\end{equation}
which is positive whenever $0 < q < 1$.
This confirms that the validation check provides strictly positive information gain over the raw dual-pixel measurement.


\section{Lennard-Jones Analogy for Cross-Modal Interaction}\label{app:lj}

The cross-modal interaction potential (Definition~\ref{def:interaction}) is inspired by the Lennard-Jones potential \citep{Jones1924}:
\begin{equation}
V_{\mathrm{LJ}}(r) = 4\varepsilon\left[\left(\frac{\sigma}{r}\right)^{12} - \left(\frac{\sigma}{r}\right)^6\right].
\end{equation}

By analogy, we can define a disparity-dependent potential:
\begin{equation}
V_{\mathrm{cross}}(D) = J\left[\left(\frac{D_0}{D + D_0}\right)^{2} - 2\left(\frac{D_0}{D + D_0}\right)\right],
\end{equation}
where $D_0$ is the characteristic disparity scale and $J$ is the coupling strength.

This potential has a minimum at $D = 0$ (the channels agree), increases for $D > 0$, and asymptotes to zero for $D \gg D_0$.
The ``attractive'' well at small disparity represents the energetic preference for channel consistency, while the repulsion at large disparity represents the cost of forcing inconsistent channels to agree.

The key difference from the molecular Lennard-Jones potential is that the variable is not spatial distance but measurement disparity---a distance in the space of categorical coordinates rather than physical space.
This mathematical analogy is structurally exact, however: in both cases, a pair of entities (atoms or measurement channels) has an equilibrium configuration (bond length or zero disparity) with a potential well controlling deviations.


%% ══════════════════════════════════════════════════════════════════════════════
%% REFERENCES
%% ══════════════════════════════════════════════════════════════════════════════
\bibliographystyle{plainnat}
\bibliography{references}

\end{document}
